% =====================================================================
% Elsevier version -- for Applied Energy / Energy and Buildings.
%
% Same body as the technical report and the IEEE build: all three
% \input ../paper/sections/*.tex, so the prose has one source. What
% this file owns is the class, Elsevier's front matter (abstract,
% keywords, structured author block, declarations) and single-column
% figure placement. Highlights are a separate file: highlights.txt.
%
% Tables: eval/paper_tables.py --out docs/paper_elsevier/tables
% =====================================================================
\documentclass[preprint,11pt]{elsarticle}

\usepackage{amsmath,amssymb}
\usepackage{graphicx}
\usepackage{booktabs}
\usepackage[final,protrusion=true,expansion=false]{microtype}
\usepackage[hidelinks]{hyperref}
% elsarticle appends a full stop to every \paragraph heading; the shared
% sections already end theirs with one. Add it only where it is missing.
\makeatletter
\def\els@aparagraph[#1]#2{\elsparagraph[#1]{#2\@addpunct.}}
\def\els@bparagraph#1{\elsparagraph*{#1\@addpunct.}}
\makeatother
\graphicspath{{../paper/figures/}}
\newcommand{\rupee}[1]{\textsf{Rs}~#1}
\usepackage{xcolor}
% Statements only the author can make. Printed in red so a submission PDF
% cannot go out with one still blank.
\newcommand{\authortodo}[1]{{\color{red}\textbf{[TODO(author): #1]}}}

% Single column, so figures sit at text width as they do in the report.
\newcommand{\paperfig}[3]{%
  \begin{figure}[t]\centering
  \includegraphics[width=0.88\textwidth]{#1}%
  \caption{#2}\label{#3}
  \end{figure}}

\journal{Applied Energy}   % chosen 2026-10-02; Energy and Buildings is the fallback

\begin{document}

\begin{frontmatter}

\title{What a Chance Constraint Does Not Promise: Measuring and
Restoring a Demand-Ceiling Guarantee on Metered Load}

% Authors and order as given 2026-10-04. All three are listed under the one
% affiliation the paper has carried; TODO(author): if an author's institution
% differs, add \affiliation[b]{...} and relabel that author [b].
\author[a]{Swetank Kumar\corref{cor1}}
\ead{dynamicsanvil@gmail.com}
\author[a]{Devansh Pathak}
\author[a]{Utkarsh Sharma}
\cortext[cor1]{Corresponding author.}
\affiliation[a]{organization={Project Samanvay / Aethergrid},
                country={India}}

\begin{abstract}
% Applied Energy caps the abstract (sources give 200 or 250 words); this is
% 199. The report and journal builds keep the longer version.
Indian commercial demand charges are set by the worst thirty-minute block in a
month, so defending a demand ceiling is a risk problem. The usual response
substitutes a high forecast quantile into the capacity constraint and calls it a
95\% guarantee. We measure what it delivers on fourteen held-out windows from an
American office and an Indian campus. Per step the bound holds, at
0.032--0.079 against a nominal 0.05; over the controller's 16-hour window the
probability of breaching somewhere is 0.39--0.82, median 0.57, eight to sixteen
times the level promised. The independence correction overstates it and is not
a bound; the valid union bound is 1 on every window. A Student-$t$ copula
predicts the realised value to a mean absolute error of 0.065, and a scenario
MILP built on it makes the violation level an operator setting. On two plants
the response is monotone and, on matched commitments, resolved within one
billing month; it is close to exact only where the per-step layer beneath it is
calibrated. Over thirty-seven supplies in eight countries, distribution-free
coverage does not degrade with aggregation, a reading this study first made,
then traced to the adaptive conformal step being stated in the series' units.
\end{abstract}

% Highlights are not in the manuscript: elsarticle prints them on a page of
% their own *before* the title page, and Elsevier takes them as a separate
% file in the submission system. They live in highlights.txt (3-5 bullets,
% at most 85 characters each; tests/test_submission.py holds them to it).

\begin{keyword}
demand charge \sep chance constraint \sep conformal prediction \sep
model predictive control \sep load forecasting \sep copula
\end{keyword}

\end{frontmatter}

\section{Introduction}

An Indian commercial electricity bill has three levers: a time-of-day energy
charge, a demand charge levied on the highest 30-minute average demand in the
month, and power-factor penalties. The second dominates the risk. It is set by
one block out of roughly $1{,}440$ in a month, so a single bad afternoon prices
all thirty days. A controller that reduces average consumption but occasionally
overshoots can cost more than one that does nothing.

The standard applied response is a chance constraint: replace the forecast mean
in the capacity constraint with a high quantile, typically the 95th, and treat
the result as a $95\%$ guarantee. Formally, with $b_t$ the uncontrollable base
load, $u_t$ the controllable allocation, $s_t$ solar generation and $D$ the
ceiling,
%
\begin{equation}
  \hat b^{\,q95}_t + u_t - \hat s^{\,q05}_t \;\le\; D
  \qquad \text{for each } t \text{ in the horizon.}
  \label{eq:marginal}
\end{equation}
%
Equation~\eqref{eq:marginal} is a set of \emph{marginal} statements. Each row
says that at that timestep, considered alone, the realisation exceeds the bound
about $5\%$ of the time. The event that costs money is different: it is
%
\begin{equation}
  \Pr\Big[\max_{t \in \mathcal{W}} \big(b_t + u_t - s_t\big) > D\Big],
  \label{eq:joint}
\end{equation}
%
a statement about a maximum over a window $\mathcal{W}$. That
\eqref{eq:marginal} does not imply \eqref{eq:joint} is elementary and well
known; joint chance constraints, their intractability, and the standard routes
around them --- Boole/Bonferroni risk allocation and the scenario approach
\citep{calafiore2006scenario,nemirovski2006convex} --- are textbook material.

\paragraph{What is not established} is how large the gap is on real building
load, and how badly each of the standard substitutes misses it. Both ends matter
practically: the first is how much a deployed controller is under-protected, and
the second is how much headroom a careful designer gives away for nothing. We
measure both, and find they are wrong in opposite directions by factors we can
state.

\paragraph{Contributions.}
\begin{enumerate}
  \item A measured bracket on horizon risk for a demand ceiling on held-out
        meter data (Section~\ref{sec:horizon}), over fourteen windows in two
        countries: claimed $0.05$, realised $0.39$--$0.82$ (median $0.57$;
        no $95\%$ lower bound under $0.28$), independence calculation
        $0.87$--$0.99$, union bound vacuous on every row, and the gap surviving
        a restatement of the per-step layer's own tuning parameter in units
        that transfer across meter sizes. The common formulation is under-protected by a factor of eight to
        sixteen, the usual correction is $1.7\times$ too conservative at the
        median, and the gap is a distribution rather than a number.
  \item A closed-loop demonstration, on two control objects in two countries
        under their own tariffs, that this risk level can be made an
        operator-set parameter rather than an emergent property of horizon
        length --- with the half of that claim which replicates separated from
        the half which does not. The response is monotone on both (rank
        correlation $0.957$ and $0.988$); it is close to exact and errs safe
        only on the plant whose per-step layer is itself calibrated, and
        repairing that layer on the other moves the dial the right way by an
        amount one billing month cannot resolve (Section~\ref{sec:dial}).
  \item A finding made and then unmade by its own audit. One protocol over
        thirty supplies in eight countries, replicated inside a single city
        across building, campus and system demand, showed conformal coverage
        degrading with aggregation. The audit of the adaptive layer at system
        scale showed why: its step is stated in the units of the series and had
        stopped moving. Re-read without that confound, split conformal fails at
        every tier alike and the adaptive layer, with its step as a fraction of
        the interval width, holds the level at every tier
        (Section~\ref{sec:calibration}). What transfers across scales is not
        the number but the fraction, and a practitioner moving the layer from a
        building to a feeder will hit exactly this.
  \item A null separating forecast skill from available flexibility as
        independent deployment axes, replicated on a panel it was not fitted to
        (Section~\ref{sec:axes}).
  \item A provenance audit of a dataset added to the study under an external
        constraint, which disqualifies that arm's most favourable result and is
        reported as such (Section~\ref{sec:china}).
\end{enumerate}

\paragraph{Where the novelty is, and where it is not.} We invented no method.
Every component --- chance-constrained MPC for buildings, joint chance
constraints and their standard treatments, quantile gradient boosting,
conformalized quantile regression, adaptive conformal inference, copula
scenario generation, decision-focused learning --- is prior work, and
Section~\ref{sec:related} places each. The novelty is empirical and
integrative: an instrument that pins everything and moves one thing, run on
metered load, with the results that argue against it reported at full size.
Section~\ref{sec:related} states the gap that survives the literature.

\section{Related work, and the gap}
\label{sec:related}

Four literatures meet at a demand ceiling, and the gap sits where they fail to
overlap.

\paragraph{Stochastic and chance-constrained MPC for buildings.}
\citet{oldewurtel2012} and \citet{oldewurtel2014} established chance-constrained
MPC for building climate control under weather uncertainty, reformulating
per-step chance constraints deterministically under Gaussian disturbance
assumptions; \citet{zhang2013scenario} replaced the distributional assumption with
sampled scenarios in the sense of \citet{calafiore2006scenario}. Joint-in-time
chance constraints --- that a trajectory stay feasible at \emph{every} step of
the horizon --- are handled in the general stochastic-MPC literature by risk
allocation under Boole's inequality \citep{ono2008risk,paulson2017joint},
which is conservative by construction and, as Section~\ref{sec:horizon}
measures, vacuous at the horizon a demand ceiling requires. A recent line
combines conformal prediction with MPC to obtain joint-in-time regions without
a distributional assumption \citep{lindemann2023safe,vogel2025conformalmpc},
evaluated on trajectory-prediction and synthetic linear-system benchmarks. The
one application of that idea to building demand response we are aware of
\citep{xu2026e2edpc} uses stage-wise --- that is, marginal --- conformal
tightening in an EnergyPlus simulation and does not report horizon-level
violation. None of these measures, on metered building load, how far a
marginally calibrated $q95$ constraint's realised horizon risk sits from its
nominal level.

\paragraph{Demand-charge management.} The closest prior work on monthly peak
limiting under load uncertainty is the Hong Kong PolyU line
\citep{xu2019adaptive,xu2020exploration}: probabilistic ANN load forecasts,
driven by probabilistic weather forecasts, feed an expected-cost threshold
resetting scheme, and Monte Carlo is used to study the robustness of the
\emph{threshold}, not to bound the probability of exceeding it.
\citet{flamm2020peak} treat demand-charge management with storage as
multi-period stochastic MPC over weather scenarios with affine decision rules
and an expected-cost objective, with a weighting for horizons that straddle
billing periods. In neither line is the exceedance probability of the ceiling
a constraint with a stated level, and neither is evaluated on an Indian tariff
or an Indian building. Field and simulation studies of HVAC demand flexibility
under MPC report peak reduction and savings without a risk statement at all.

\paragraph{Conformal prediction for multi-step forecasts.} Split conformal
\citep{vovk2005algorithmic} and CQR \citep{romano2019cqr} give per-step
marginal coverage. \citet{stankeviciute2021conformal} obtain joint coverage
over a horizon by Bonferroni correction; \citet{sun2024copula} replace
Bonferroni with a copula over per-step conformity scores, which is the same
modelling move as Section~\ref{sec:dial} but deployed as a \emph{predictor}
rather than inside an optimiser, and requiring a second calibration set.
\citet{lin2022temporal} study conformal intervals under temporal dependence.
\citet{zaffran2022adaptive} and \citet{wang2024acmcp} extend adaptive conformal
inference \citep{gibbs2021adaptive} to non-stationary and multi-step settings,
targeting marginal coverage; the latter establishes that optimal $h$-step
errors are serially correlated to lag $h-1$, which is the dependence our
copula fits. This paper adds to that line a measurement rather than a method:
the realised joint exceedance on real load errors, placed against both the
Bonferroni bound and the independence figure, with the copula's accuracy for
that quantity reported.

\paragraph{Decision-aware forecasting and calibration.} Decision-focused
learning \citep{elmachtoub2022spo,donti2017task,berthet2020perturbed} trains
the forecast on the downstream decision's cost; \citet{stratigakos2026decision}
calibrate multivariate prediction sets to the \emph{reliability of the
downstream decision} rather than to coverage, for robust reserve scheduling.
That is complementary to Section~\ref{sec:dial}, which reaches the same object
--- an operator-set violation level on a committed ceiling --- from the
optimiser side, and reports its exactness and resolution floor.

\paragraph{Data.} Probabilistic load forecasting is mature at system scale
\citep{hong2016review,hong2016gefcom} and increasingly benchmarked at building
scale \citep{miller2020bdg2,emami2023buildingsbench,nweye2023citylearn}. Aggregation is known to
improve point accuracy by a scaling law \citep{sevlian2018scaling}; whether
\emph{distribution-free coverage} holds across aggregation levels has not, to
our knowledge, been reported. For India, \citet{rashid2019iblend} publish
I-BLEND: 52 months of one-minute mains data, with power factor, for seven
buildings on a Delhi campus. It is the only open Indian building-level load
corpus at the resolution Indian demand charges are assessed on, and it has
been used for point forecasting rather than for demand-charge control.

\paragraph{The gap.} We find no published work that (i) measures, on metered
building load, the horizon-level breach probability a marginally calibrated
$q95$ chance constraint actually delivers against its nominal level, bracketed
by the valid Boole bound above and the independence figure that is usually
misquoted as one; (ii) restores that level as a settable parameter inside a
demand-charge MILP and reports its exactness and resolution floor; (iii) tests
distribution-free coverage across building and system scale under one
protocol; or (iv) does any of this on real Indian commercial buildings at
native 15-minute resolution under an Indian tariff. Items (i)--(iii) are this
paper. Item (iv) is the work Section~\ref{sec:india-fix} opens, and the data
for it exists.

\section{The instrument}
\label{sec:instrument}

Six components, each independently swappable, with the bill engine as the single
source of truth for every reported figure:
%
\begin{align*}
\text{tariff} &\;\rightarrow\; \text{quantile forecast}\\
  &\;\rightarrow\; \text{conformal calibration}\\
  &\;\rightarrow\; \text{chance-constrained MILP}\\
  &\;\rightarrow\; \text{RC simulator}
    \;\rightarrow\; \text{bill engine}.
\end{align*}
%
The forecaster is a quantile LightGBM \citep{ke2017lightgbm} at five levels
over a 64-step (16-hour) horizon at 15-minute resolution, wrapped in split and
adaptive conformal calibration
\citep{angelopoulos2023gentle,gibbs2021adaptive}. The controller is a
deterministic MILP solved with HiGHS; no
language model touches a money decision. The evaluation discipline is that a
number reaches a table only if a script produced it, and a single entry point
regenerates every figure and every table in this paper from a clean checkout.

This matters for the measurements because each one requires pinning everything
and moving one thing. Section~\ref{sec:horizon} pins the plant and varies only
the risk formulation. Section~\ref{sec:calibration} pins the protocol and varies
only the series. Section~\ref{sec:axes} pins the model and varies the site.
Section~\ref{sec:china} pins all three and varies the country, which is the one
we did not choose.

\section{Data}
\label{sec:data}

Six public sources, all fetched by script rather than vendored, and named in
full because the benchmark's credibility rests on them.

\paragraph{Building Data Genome Project 2} \citep{miller2020bdg2}. $1{,}636$
metered non-residential buildings across nineteen sites, of which $1{,}578$
carry an electricity meter, hourly with matched weather, two years. It is the
only harmonised multi-country building corpus in existence, which is why it
carries the controlled comparison: the same physical quantity, the same metering
convention, the same period. It contains \emph{no Indian buildings}, which is
why the next source is in the study.

\paragraph{I-BLEND} \citep{rashid2019iblend}. One-minute mains power and power
factor for seven buildings on the IIIT-Delhi campus --- academic, lecture,
library, facilities, dining and two dormitories --- and for the three
transformers that feed the whole campus, from August 2013 to December 2017,
CC0. The billing blocks are formed by averaging sixty one-minute readings, not
by interpolating up from hourly data, so this is the only building-level
source in the study at the cadence Indian demand charges are assessed on. Its
documentation states, and the load--temperature signature of each meter
confirms, that every building except Facilities is served by a central chiller
plant that is \emph{not} on the building's own meter: at building level the
corpus is therefore a forecasting and calibration benchmark rather than a
control one, and the selection rule's chilled-water clause records it as such.
The campus feed is the exception, and it is the more important object: an
Indian campus is billed as a single high-tension consumer with one contract
demand and one monthly demand charge, so the transformer total --- median
$199$\,kW, peak $842$\,kW, chiller included --- is what the tariff actually
sees. Meter uptime varies by feed; each series' admissible test Junes and the
training months behind each are recorded by \texttt{data/iblend.py} and
printed in the table.

\paragraph{Delhi State Load Despatch Centre.} Fifteen-minute metered demand for
the city of Delhi, disaggregated across five distribution companies (BRPL, BYPL,
NDPL, NDMC and MES) plus the total, two complete years at full coverage. With
I-BLEND this gives three levels of aggregation --- building, campus, city --- in
one city, at native resolution, under one weather feed. Every non-Indian series
in the study is interpolated up from hourly.

\paragraph{ENTSO-E transmission data, via Open Power System Data.} Hourly
national load for the United Kingdom, Ireland, Germany, France and Spain, drawn
from the European transmission system operators' transparency platform.

\paragraph{Chinese provincial demand} \citep{wu2023china}. Hourly load for the
31 provinces of mainland China, 2018, from which six are drawn. This source was
added after the study was otherwise complete, and it is the only tier-2 series
that is constructed rather than metered. Section~\ref{sec:china} audits it.

\paragraph{ERA5 reanalysis, via the Open-Meteo archive.} Hourly temperature, dew
point and cloud cover at each location, used identically for every series.

\paragraph{Windows.} The building tier and the European national tier share
2016--2017. The Delhi archive covers 2011--2012, so the same protocol is shifted
back five years with its shape preserved. The Chinese series is a single
calendar year, 2018. In every row of the study the test month is June, and the
headline metric is skill against each series' own seasonal-naive baseline: a
ratio computed inside one series' own data and its own period, which is what
makes rows with different units and different years comparable at all. Absolute
losses are not comparable across groups and are never compared.

\paragraph{Selection.} One rule, stated in advance and applied uniformly: no
district chilled-water meter, $>$$97\%$ coverage, median load 40--800\,kW,
positive load--temperature correlation; ties within a (site, usage) cell broken
by highest $p99/\text{median}$. Thirty supplies pass. Two facts about the rule
were discovered by applying it rather than anticipated when writing it, and both
are reported in Section~\ref{sec:calibration}.

\section{The horizon gap}
\label{sec:horizon}

\begin{table}[t]
\centering
\footnotesize
\setlength{\tabcolsep}{3.5pt}
\resizebox{\linewidth}{!}{%
\begin{tabular}{rrrrcrrr}
\toprule
$H$ & hours & per-step $\hat\alpha$ & realised & 95\% CI & independent & Boole $H\hat\alpha$ & copula \\
\midrule
1 & 0.25 & 0.060 & \textbf{0.060} & [0.03, 0.09] & 0.060 & 0.060 & 0.049 \\
4 & 1.00 & 0.060 & \textbf{0.091} & [0.05, 0.13] & 0.219 & 0.240 & 0.100 \\
8 & 2.00 & 0.060 & \textbf{0.126} & [0.08, 0.18] & 0.389 & 0.478 & 0.148 \\
16 & 4.00 & 0.059 & \textbf{0.182} & [0.11, 0.25] & 0.625 & 0.952 & 0.225 \\
32 & 8.00 & 0.061 & \textbf{0.275} & [0.18, 0.38] & 0.867 & 1.000 & 0.342 \\
64 & 16.00 & 0.062 & \textbf{0.405} & [0.28, 0.53] & 0.983 & 1.000 & 0.497 \\
\bottomrule
\end{tabular}
}
\caption{Marginal versus horizon-level exceedance of the demand ceiling, held-out June 2017, 2,817 forecast origins over 31 days. The per-step rate $\hat\alpha$ is well calibrated at every horizon against a nominal $0.05$. The realised probability of breaching \emph{somewhere} in the window is not that number: at $H=64$ it is between $6$ and $11$ times the nominal level (day-block bootstrap, 95\% interval; adjacent origins share 63 of 64 steps, so the day is the unit of replication). Neither substitute is usable: the independence calculation $1-(1-\hat\alpha)^H$ overstates the risk by $2.4\times$ because load errors are strongly autocorrelated, and the Boole/Bonferroni union bound $\min(1, H\hat\alpha)$ -- the only one of the two that is a valid bound -- saturates at $1$ by $H=17$ and is vacuous over the controller's actual horizon.}
\label{tab:horizon}
\end{table}


\paperfig{fig_horizon.pdf}{The bracket. Marginal calibration is a statement about each row of
\eqref{eq:marginal} separately; it constrains the probability of breaching
somewhere in the window only to the shaded region, which at $H=64$ spans
$0.05$ to $0.983$. The per-step rate (bottom) is flat and well calibrated at
every horizon, so nothing here is a broken forecaster. The band around the
realised curve is the $95\%$ day-block bootstrap interval. Held-out June 2017,
one 13{,}759\,m$^2$ office, $2{,}817$ forecast origins.}{fig:horizon}

Table~\ref{tab:horizon} and Figure~\ref{fig:horizon} report, for a
13{,}759\,m$^2$ office on held-out June 2017 data over $2{,}817$ forecast
origins, the per-step exceedance rate of the ceiling alongside the probability
that the ceiling is breached at least once within a window of $H$ steps.

Three readings of the last row are simultaneously true and together are the
result.

\paragraph{The forecaster is not the problem.} The per-step quantile is well
calibrated at every horizon: $\hat\alpha$ holds between $0.059$ and $0.062$
against a nominal $0.05$. Marginal calibration delivers exactly what it
promised, which is why this is a result about a formulation rather than a bug
report about a model.

\paragraph{The formulation is under-protected by a factor of eight.} The
probability of breaching somewhere in the 16-hour window is $0.405$: $6.5\times$
the per-step rate the forecaster actually delivers, and $8.1\times$ the $0.05$
that \eqref{eq:marginal} appears to promise. That figure carries sampling
error, and the naive kind would understate it: adjacent origins share 63 of
their 64 steps, so the unit of replication is the day rather than the origin,
and a day-block bootstrap puts the $95\%$ interval at $0.28$--$0.53$, between
$5.7\times$ and $10.7\times$ nominal. A controller deployed on
\eqref{eq:marginal} and described to an operator as a $95\%$ guarantee is
running at roughly $60\%$, and at no better than $72\%$.

\paragraph{Neither standard substitute recovers it, and they fail in opposite
directions.} This is the correction to the usual telling, and it matters. The
quantity commonly quoted, $1-(1-\alpha)^H$, is the \emph{independence}
calculation, not a bound: it assumes the exceedance events are independent,
which is the one thing load errors are not. It gives $0.983$, which is
$2.4\times$ the realised $0.405$, because the fitted error process is heavily
autocorrelated --- lag-1 correlation $0.871$, $0.578$ at lag 4, decaying to
$0.085$ by lag 32 --- so exceedances cluster rather than spread. The quantity
that \emph{is} a valid bound, Boole/Bonferroni's $\min(1, H\hat\alpha)$, makes
no dependence assumption and pays for it: at $\hat\alpha \approx 0.06$ it
reaches $1$ by $H=17$, four hours into a sixteen-hour horizon, after which it
says nothing at all.

The practical consequence runs in both directions. A controller using
\eqref{eq:marginal} is materially under-protected. A controller sized off the
independence figure is over-protected by more than a factor of two and pays for
it in usable headroom, which is the quantity the entire economic argument rests
on. A controller that insists on a valid bound over this horizon has no bound to
use.

\begin{table}[t]
\centering
\footnotesize
\setlength{\tabcolsep}{3.5pt}
\resizebox{\linewidth}{!}{%
\begin{tabular}{llrrcrrrr}
\toprule
Series & June & per-step $\hat\alpha$ & realised $H{=}64$ & 95\% CI & independent & Boole & copula & lag-1 $\rho$ \\
\midrule
Phoenix office (BDG2) & 2017 & 0.062 & \textbf{0.405} & [0.28, 0.53] & 0.983 & 1.00 & 0.497 & 0.87 \\
\midrule
IIIT-Delhi Academic & 2015 & 0.049 & \textbf{0.393} & [0.28, 0.50] & 0.959 & 1.00 & 0.477 & 0.77 \\
IIIT-Delhi Academic & 2016 & 0.050 & \textbf{0.636} & [0.50, 0.75] & 0.963 & 1.00 & 0.580 & 0.72 \\
IIIT-Delhi Academic & 2017 & 0.048 & \textbf{0.422} & [0.30, 0.54] & 0.956 & 1.00 & 0.595 & 0.71 \\
IIIT-Delhi Boys & 2014 & 0.032 & \textbf{0.501} & [0.38, 0.61] & 0.874 & 1.00 & 0.554 & 0.72 \\
IIIT-Delhi Campus & 2017 & 0.079 & \textbf{0.432} & [0.29, 0.58] & 0.995 & 1.00 & 0.482 & 0.85 \\
IIIT-Delhi Dining & 2017 & 0.053 & \textbf{0.714} & [0.61, 0.81] & 0.970 & 1.00 & 0.744 & 0.72 \\
IIIT-Delhi Facilities & 2015 & 0.052 & \textbf{0.639} & [0.55, 0.71] & 0.967 & 1.00 & 0.705 & 0.63 \\
IIIT-Delhi Facilities & 2016 & 0.054 & \textbf{0.621} & [0.52, 0.72] & 0.972 & 1.00 & 0.633 & 0.66 \\
IIIT-Delhi Facilities & 2017 & 0.048 & \textbf{0.739} & [0.62, 0.85] & 0.958 & 1.00 & 0.753 & 0.61 \\
IIIT-Delhi Girls & 2014 & 0.048 & \textbf{0.693} & [0.59, 0.78] & 0.957 & 1.00 & 0.589 & 0.63 \\
IIIT-Delhi Girls & 2017 & 0.046 & \textbf{0.818} & [0.74, 0.89] & 0.951 & 1.00 & 0.739 & 0.46 \\
IIIT-Delhi Library & 2014 & 0.036 & \textbf{0.480} & [0.36, 0.59] & 0.902 & 1.00 & 0.567 & 0.73 \\
IIIT-Delhi Library & 2016 & 0.042 & \textbf{0.509} & [0.41, 0.61] & 0.935 & 1.00 & 0.504 & 0.80 \\
\bottomrule
\end{tabular}
}
\caption{The bracket of Table~\ref{tab:horizon}, replicated over 14 windows: the Phoenix office of the original measurement and 13 Indian windows at native 15-minute resolution. The per-step rate is calibrated throughout ($0.032$--$0.079$ against a nominal $0.05$). The realised probability of breaching somewhere in the 16-hour window runs from $0.393$ to $0.818$, median $0.565$ --- between $8$ and $16$ times the level the constraint appears to promise. Intervals are day-block bootstrap, the day being the unit of replication because adjacent origins share 63 of 64 steps. No window's lower bound falls under $0.28$, $5.7\times$ nominal. The spread across windows is not sampling noise: 32 of 91 pairs of intervals are disjoint, the lowest and highest rows among them, and the same building returns different values in different years while its per-step rate does not move. A 16-hour window opened late in one day runs into the next, so two-day blocks are also reported: they widen the intervals by 13\%, and the floor becomes $0.25$ ($5.0\times$) with 20 pairs disjoint. The independence figure overstates realised risk by $1.66\times$ (median) and lies inside the realised interval on 0 of 14 windows; the Boole bound is $1$ on every row. The copula predicts the realised value to a mean absolute error of $0.065$ against $0.382$ for independence, and lies inside the interval on 12 of 14 windows.}
\label{tab:horizon-panel}
\end{table}


\paperfig{fig_horizon_panel.pdf}{The bracket, one row per window, sorted by the realised value. The
per-step rate sits on the nominal line on every row; the realised horizon
exceedance does not sit anywhere in particular, and the same meter returns
different values in different Junes. The independence figure is near $1$
throughout. Phoenix is the original measurement; the rest are IIIT-Delhi at
native 15-minute resolution.}{fig:horizon-panel}

\paragraph{The bracket replicates, and it is a distribution.} The measurement
above is one building in one month. Table~\ref{tab:horizon-panel} and
Figure~\ref{fig:horizon-panel} repeat it, with nothing changed but the series,
over thirteen windows on the IIIT-Delhi campus of Section~\ref{sec:data} ---
six buildings and the campus feed, up to three Junes each, every one at native
15-minute resolution rather than interpolated from hourly. Three things hold on
every row. The per-step rate is calibrated, $0.032$--$0.079$ against a nominal
$0.05$. The Boole bound is $1$. And the independence figure sits between
$0.87$ and $0.99$, overstating the realised risk by $1.66\times$ at the median.
What does not hold is the number itself. The realised probability of breaching
somewhere in the window runs from $0.393$ to $0.818$, median $0.565$ ---
between eight and sixteen times the level the constraint appears to promise ---
and it is not a constant of the formulation: the Academic building returns
$0.393$, $0.636$ and $0.422$ in three successive Junes while its per-step rate
stays within $0.002$ of $0.05$. The Phoenix figure of $0.405$ was at the mild
end of the range. The intervals say the spread is not sampling noise: no
window's lower bound falls under $0.28$, $5.7\times$ nominal, and 32 of the 91
pairs of intervals are disjoint, the lowest and highest rows among them. Two-day
blocks, which allow for a window opened late in one day running into the next,
widen the intervals by $13\%$ and leave 20 pairs disjoint. The gap is a
property of how exceedances cluster in a given month, which marginal calibration says nothing about, and the practical reading
is that an operator cannot even be told a single corrected number: the
correction is itself a random variable with a wide law.

The copula is the one estimate that tracks it. Across the fourteen windows it
predicts the realised value to a mean absolute error of $0.065$, against
$0.382$ for the independence figure, and lies inside the realised value's
$95\%$ interval on twelve of the fourteen; the independence figure lies inside
none. It is not exact --- it misses the two most extreme rows by $0.11$ and $0.13$ --- and
Section~\ref{sec:dial} states its limits. But it is the only quantity in the
table that moves with the truth, and that is the property the dial in the next
section depends on.

\begin{table}[t]
\centering
\footnotesize
\setlength{\tabcolsep}{3.5pt}
\resizebox{\linewidth}{!}{%
\begin{tabular}{llrrrrcr}
\toprule
Series & June & $\hat\alpha$ abs. & $\hat\alpha$ rel. & realised abs. & realised rel. & 95\% CI & copula \\
\midrule
Phoenix office (BDG2) & 2017 & 0.061 & 0.072 & 0.426 & \textbf{0.425} & [0.29, 0.56] & 0.452 \\
\midrule
IIIT-Delhi Academic & 2015 & 0.049 & 0.047 & 0.393 & \textbf{0.342} & [0.24, 0.45] & 0.428 \\
IIIT-Delhi Academic & 2016 & 0.050 & 0.055 & 0.636 & \textbf{0.572} & [0.43, 0.71] & 0.481 \\
IIIT-Delhi Academic & 2017 & 0.048 & 0.049 & 0.422 & \textbf{0.351} & [0.22, 0.49] & 0.441 \\
IIIT-Delhi Boys & 2014 & 0.032 & 0.025 & 0.501 & \textbf{0.325} & [0.22, 0.43] & 0.420 \\
IIIT-Delhi Campus & 2017 & 0.079 & 0.069 & 0.432 & \textbf{0.438} & [0.30, 0.57] & 0.528 \\
IIIT-Delhi Dining & 2017 & 0.053 & 0.054 & 0.714 & \textbf{0.660} & [0.55, 0.76] & 0.656 \\
IIIT-Delhi Facilities & 2015 & 0.052 & 0.074 & 0.639 & \textbf{0.603} & [0.51, 0.69] & 0.602 \\
IIIT-Delhi Facilities & 2016 & 0.054 & 0.082 & 0.621 & \textbf{0.609} & [0.49, 0.72] & 0.546 \\
IIIT-Delhi Facilities & 2017 & 0.048 & 0.053 & 0.739 & \textbf{0.516} & [0.37, 0.65] & 0.614 \\
IIIT-Delhi Girls & 2014 & 0.048 & 0.034 & 0.693 & \textbf{0.442} & [0.31, 0.58] & 0.378 \\
IIIT-Delhi Girls & 2017 & 0.046 & 0.042 & 0.818 & \textbf{0.632} & [0.51, 0.75] & 0.450 \\
IIIT-Delhi Library & 2014 & 0.036 & 0.033 & 0.480 & \textbf{0.254} & [0.15, 0.37] & 0.399 \\
IIIT-Delhi Library & 2016 & 0.042 & 0.026 & 0.509 & \textbf{0.307} & [0.20, 0.40] & 0.378 \\
\bottomrule
\end{tabular}
}
\caption{The horizon rows of Table~\ref{tab:horizon-panel} with the per-step layer's adaptive step stated as a fraction of the interval width ($\kappa=0.006$, the width read on the validation block) instead of as an absolute number in the units of each series. These meters run from intervals of a few kilowatts to several hundred, so the absolute step was a different layer on every row --- $\kappa=0.0015$ to $\kappa=0.0691$ across these windows, so on all but the campus feed it was faster than the step the building audit settled on, not slower. Under the relative step the per-step rate spans $0.025$--$0.082$ against $0.032$--$0.079$ before, and its mean absolute distance from the nominal $0.05$ moves from $0.0073$ to $0.0144$: a slower layer tracks the nominal rate less tightly, which is why the per-step column of Table~\ref{tab:horizon-panel} is the one reported there. The horizon gap survives and shrinks. Realised 16-hour exceedance runs $0.254$--$0.660$ (median $0.440$, $9\times$ nominal) against $0.393$--$0.818$ (median $0.565$) before; the paired change is $-0.111$ on average, and the absolute-step value lies inside the relative-step row's own $95\%$ interval on 8 of 14 windows, so on 6 the two are separated by more than sampling error. The magnitude of the gap is therefore not independent of how the per-step layer is tuned --- a faster layer produces a noisier bound and more nearly independent exceedances across the horizon --- but no choice of step makes it go away: the lowest lower bound anywhere in this table is $0.15$, $3.0\times$ the level the constraint appears to promise. The copula lies inside the realised interval on 12 of 14 windows. Fox is run through this trainer at both steps, because its row in Table~\ref{tab:horizon-panel} comes from the benchmark harness and would not otherwise be a paired comparison.}
\label{tab:horizon-kappa}
\end{table}


\paragraph{The gap is not an artefact of how the per-step layer was tuned, but
its size is not independent of it either.} Every row above was calibrated with
an adaptive layer whose step is a number in the units of its own series, and
these meters run from intervals of five kilowatts to two hundred and twenty.
Section~\ref{sec:calibration} shows what that did to the calibration table;
the same objection applies here, because a per-step rate is what the first
column reports. Table~\ref{tab:horizon-kappa} repeats all fourteen windows
with the step restated as a fraction of the interval width, $\kappa=0.006$,
the width read on the validation block. Two things follow, and only one of them
was expected. The gap survives: realised 16-hour exceedance runs $0.254$ to
$0.660$, median $0.440$, nine times the level the constraint appears to
promise, and the lowest bound anywhere in the table is $0.15$, three times
nominal. But the gap also shrinks, by $0.111$ on average, and on six of the
fourteen windows the two measurements are separated by more than sampling
error. The mechanism is visible in the first two columns: on these small meters
the absolute step was $\kappa=0.02$ to $\kappa=0.07$, a \emph{faster} layer
than the relative one and not a slower one, and a fast layer produces a bound
that jumps after each exceedance and decays between them. That bound is
noisier, its exceedances are closer to independent across the horizon, and the
probability that at least one occurs is correspondingly higher. A slower layer
gives a smoother bound, more clustered exceedances, and a smaller joint
probability --- at the cost of tracking the nominal per-step rate less
tightly, which is why the panel above, whose per-step column is the better
calibrated of the two ($0.0073$ against $0.0144$ mean absolute distance from
nominal), is the one reported as the measurement.

The honest statement is therefore narrower than ``the gap is $8$ to $16$
times'' and survives being narrowed. Under both calibrations of the per-step
layer, on every one of fourteen windows in two countries, a marginally
calibrated $q95$ constraint delivers a horizon-level breach probability
between three and sixteen times its nominal level, and which end of that range
a particular deployment lands on depends on a tuning choice nobody making it
would expect to be choosing a risk level.

\section{Restoring the guarantee as a dial}
\label{sec:dial}

\begin{table}[t]
\centering
\footnotesize
\setlength{\tabcolsep}{3.5pt}
\resizebox{\linewidth}{!}{%
\begin{tabular}{llrrcrrrrr}
\toprule
Target & Controller & $\varepsilon$ & Commit viol.\ (gap) & 95\% CI & vs target & Breaches & Peak kVA & Bill Rs & Solve ms \\
\midrule
tight & marginal $q95$ & -- & \textbf{0.106} & [0.03, 0.20] & 0.049 & 4 & 459.0 & 1,138,193 & 6 \\
tight & scenario & 0.05 & \textbf{0.097} (+0.047) & [0.04, 0.17] & 0.097 & 9 & 464.6 & 1,141,532 & 52 \\
tight & scenario & 0.10 & \textbf{0.114} (+0.014) & [0.04, 0.21] & 0.096 & 9 & 466.8 & 1,142,186 & 52 \\
tight & scenario & 0.20 & \textbf{0.156} (-0.044) & [0.07, 0.25] & 0.073 & 10 & 469.2 & 1,142,647 & 55 \\
tight & scenario & 0.35 & \textbf{0.219} (-0.131) & [0.13, 0.33] & 0.077 & 13 & 467.5 & 1,141,036 & 55 \\
nominal & marginal $q95$ & -- & \textbf{0.103} & [0.03, 0.19] & 0.000 & 0 & 481.9 & 1,149,580 & 6 \\
nominal & scenario & 0.05 & \textbf{0.078} (+0.028) & [0.02, 0.15] & 0.000 & 0 & 480.5 & 1,148,576 & 63 \\
nominal & scenario & 0.10 & \textbf{0.097} (-0.003) & [0.03, 0.18] & 0.000 & 0 & 486.0 & 1,151,865 & 63 \\
nominal & scenario & 0.20 & \textbf{0.150} (-0.050) & [0.07, 0.24] & 0.000 & 0 & 480.1 & 1,148,071 & 63 \\
nominal & scenario & 0.35 & \textbf{0.208} (-0.142) & [0.12, 0.31] & 0.000 & 0 & 484.0 & 1,150,364 & 55 \\
\bottomrule
\end{tabular}
}
\caption{Closed loop over one billing month on a 13{,}759\,m$^2$ office in Phoenix under the TNERC tariff. \textbf{Commit violation} is the acceptance metric: the fraction of horizons in which realised load cleared the ceiling the optimiser committed to, which is what $\varepsilon$ is a statement about. The against-target column is the business metric and is not what the chance constraint promises, since the committed peak is a decision variable. Rank correlation 0.957, mean absolute gap 0.057, conservative at 5 of 8 levels; resolution floor $1/S=0.025$. Intervals are a day-block bootstrap over the sampled commitments, which are overlapping 16-hour windows and so are not the independent draws a binomial interval would assume; the requested $\varepsilon$ lies inside its own interval at 6 of 8 levels, and the mean half-width is $0.081$, which is the resolution one billing month buys.}
\label{tab:acceptance}
\end{table}


We fit a Student-$t$ copula over the horizon and sample joint trajectories
--- the marginals-to-scenarios construction of \citet{pinson2009scenarios},
with a heavier-tailed dependence family --- which a scenario MILP then
constrains with a CVaR surrogate \citep{rockafellar2000cvar,nemirovski2006convex}
on the probability mass permitted to violate. The copula reproduces the calibrated marginals to
within $4.6\%$ of mean interval width at $4{,}000$ paths per origin, so it adds
dependence and disturbs nothing upstream, and it predicts horizon exceedance far
closer than either substitute: $0.497$ against a realised $0.405$ (interval
$0.28$--$0.53$) at $H=64$,
where \eqref{eq:marginal} implies $0.05$ and the independence figure is $0.983$.

Table~\ref{tab:acceptance} is a closed-loop run over a real billing month at two
demand targets and four risk levels. The acceptance metric is violation of the
ceiling \emph{the optimiser itself committed to at each origin}, because that is
the event $\varepsilon$ is a statement about. The operator's monthly demand
target is a business metric the chance constraint never promised: the committed
peak is a decision variable, and an optimiser short of capacity may lift it, pay
the demand charge, and honour its constraint exactly. Grading a chance
constraint against a number it never promised is the usual way this test is made
meaningless, so both columns are shown and only one is the acceptance criterion.

The result is that $\varepsilon$ becomes a dial. Realised commit violation
tracks the requested level with rank correlation $0.957$ across the sweep, mean
absolute gap $0.057$, and lands conservative at five of eight levels. The
formulation errs safe by construction: the CVaR surrogate admits a strict subset
of the chance-constrained feasible set, so it defends the committed ceiling
harder than asked. Those are point estimates over one month, and the
commitments behind them are overlapping sixteen-hour windows, so each rate
carries a day-block interval; the mean half-width is $0.081$. The requested
$\varepsilon$ lies inside its own interval at six of eight levels, and the two
it does not are both $\varepsilon=0.35$, on the safe side: on this plant the
dial is never measurably unsafe.

\begin{table}[t]
\centering
\footnotesize
\setlength{\tabcolsep}{3.5pt}
\resizebox{\linewidth}{!}{%
\begin{tabular}{llrrcrrrrr}
\toprule
Target & Controller & $\varepsilon$ & Commit viol.\ (gap) & 95\% CI & vs target & Breaches & Peak kVA & Bill Rs & Solve ms \\
\midrule
tight & marginal $q95$ & -- & \textbf{0.297} & [0.18, 0.42] & 0.101 & 13 & 890.9 & 2,837,240 & 6 \\
tight & scenario & 0.05 & \textbf{0.178} (+0.128) & [0.08, 0.28] & 0.101 & 14 & 885.5 & 2,833,912 & 51 \\
tight & scenario & 0.10 & \textbf{0.233} (+0.133) & [0.13, 0.35] & 0.124 & 15 & 895.3 & 2,835,693 & 61 \\
tight & scenario & 0.20 & \textbf{0.300} (+0.100) & [0.19, 0.42] & 0.123 & 14 & 903.8 & 2,838,637 & 62 \\
tight & scenario & 0.35 & \textbf{0.392} (+0.042) & [0.26, 0.52] & 0.075 & 9 & 915.4 & 2,842,287 & 63 \\
nominal & marginal $q95$ & -- & \textbf{0.297} & [0.18, 0.42] & 0.000 & 0 & 914.0 & 2,843,724 & 6 \\
nominal & scenario & 0.05 & \textbf{0.172} (+0.122) & [0.08, 0.27] & 0.023 & 1 & 945.5 & 2,849,643 & 61 \\
nominal & scenario & 0.10 & \textbf{0.236} (+0.136) & [0.13, 0.36] & 0.023 & 1 & 932.7 & 2,846,024 & 63 \\
nominal & scenario & 0.20 & \textbf{0.300} (+0.100) & [0.19, 0.42] & 0.023 & 1 & 931.3 & 2,846,298 & 81 \\
nominal & scenario & 0.35 & \textbf{0.392} (+0.042) & [0.26, 0.52] & 0.023 & 1 & 938.9 & 2,848,913 & 59 \\
\bottomrule
\end{tabular}
}
\caption{Closed loop over one billing month on the IIIT-Delhi campus feed under the DERC tariff, the billed consumer of Section~\ref{sec:data}. \textbf{Commit violation} is the acceptance metric: the fraction of horizons in which realised load cleared the ceiling the optimiser committed to, which is what $\varepsilon$ is a statement about. The against-target column is the business metric and is not what the chance constraint promises, since the committed peak is a decision variable. Rank correlation 0.988, mean absolute gap 0.100, conservative at 0 of 8 levels; resolution floor $1/S=0.025$. Intervals are a day-block bootstrap over the sampled commitments, which are overlapping 16-hour windows and so are not the independent draws a binomial interval would assume; the requested $\varepsilon$ lies inside its own interval at 4 of 8 levels, and the mean half-width is $0.113$, which is the resolution one billing month buys.}
\label{tab:acceptance-campus}
\end{table}


\paragraph{A second control object, and the half of this that replicates.}
Table~\ref{tab:acceptance-campus} repeats the sweep on the IIIT-Delhi campus
feed --- another country, another tariff, another plant, and the consumer a
Delhi utility actually bills --- with nothing tuned for it. Two things happen,
and only one of them is the result above.

The dial's \emph{monotonicity} replicates, and more cleanly than on the office:
rank correlation $0.988$, and realised commit violation rises through $0.172$,
$0.236$, $0.300$, $0.392$ as $\varepsilon$ is set to $0.05$, $0.10$, $0.20$,
$0.35$. An operator turning the dial gets a monotone response on both plants.

The dial's \emph{calibration} does not replicate. Mean absolute gap is $0.100$
against $0.057$, and the sweep is conservative at \emph{none} of its eight
levels rather than five: on this feed the formulation delivers more risk than
it was asked for at every setting. The intervals sharpen where that matters:
they exclude the requested level at $\varepsilon=0.05$ and $0.10$ on both
targets --- the two tightest settings, the ones an operator would actually use
--- and contain it at $0.20$ and $0.35$, four of eight in all. We do not think this is the risk
formulation failing, and the reason is one row of
Table~\ref{tab:horizon-panel}. The campus feed is the single window in the
study whose per-step layer \emph{under}-covers --- $\hat\alpha = 0.079$
against a nominal $0.05$, the only figure above the line --- and $\varepsilon$
is a statement built on top of that layer. A joint guarantee cannot be more
exact than the marginals it composes, and here the marginals are optimistic by
half again.

That is a testable claim, so we tested it. Restating the adaptive step in the
interval's units --- which Section~\ref{sec:calibration} shows is the
defensible way to set it --- moves the campus per-step rate from $0.079$ to
$0.069$, and the whole sweep was rerun on that calibration with nothing else
changed. The dial's exactness improves in the direction the explanation
predicts, and by an amount this experiment cannot separate from noise. Mean
absolute gap falls from $0.100$ to $0.083$ and realised violation at
$\varepsilon=0.05$ from $0.178$ to $0.158$, both toward the $0.057$ and $0.097$
of the well-calibrated plant; rank correlation is unchanged at $0.988$ and the
sweep is still conservative at none of its eight levels. But each of those
rates is a fraction of the commitments sampled over one month, and the
day-block interval on it is wider than the movement --- a mean half-width of
$0.113$ on the campus sweep against a shift of $0.018$ --- so what this run
establishes is a direction and not a size. The claim the paper makes is
therefore the weak one: a per-step layer that under-covers is a plausible and
partially-supported cause of a dial that over-delivers risk, the repair moves
the dial the right way, and the remaining distance to the calibrated plant is
neither shown to close nor shown to persist. Separating those would need more
billing months than one.

The honest summary of Contribution~2 is therefore narrower than one control
object suggested. On both plants $\varepsilon$ converts a risk level that was
an emergent property of horizon length into one an operator sets and gets a
monotone response from. On the plant whose marginals are calibrated it is also
roughly exact and errs safe; on the plant whose marginals are not, it is
monotone and optimistic, and repairing those marginals moves it the right way
by an amount one month cannot resolve. That is a useful thing for a deployment to know, and
it says the marginal calibration layer is not a preliminary to the guarantee
but part of it.

\paragraph{What it does not do, stated in full.} It is not exact. At the
tightest setting, $\varepsilon=0.05$ delivers $0.097$ at the tight target,
nearly twice the requested risk and on the unsafe side --- though one month's
interval on it, $0.04$--$0.17$, cannot tell that miss from none. The dial has
two resolution limits. One is statistical: the intervals are wider than the
spacing between adjacent levels, so the monotone ordering above is a property
of the point estimates, and resolving neighbouring settings one from another
would take more billing months, not more scenarios. The other is a hard floor
at $1/S = 0.025$: with $S$ scenarios the finest expressible
risk level is one scenario's worth of probability mass, below which the CVaR
surrogate collapses onto the sample maximum and successive $\varepsilon$ cease
to be distinguishable. That is a property of sample average approximation rather
than a tuning artefact, and a finer dial costs scenarios and solve time.

\begin{table}[t]
\centering
\footnotesize
\setlength{\tabcolsep}{3.5pt}
\resizebox{\linewidth}{!}{%
\begin{tabular}{llrcrrrrr}
\toprule
Series & June & realised & 95\% CI & $\nu_{\mathrm{tail}}$ & copula & $\nu_{\mathrm{lik}}$ & copula & copula, $\nu{=}7$ \\
\midrule
Phoenix office (BDG2) & 2017 & \textbf{0.405} & [0.28, 0.53] & 7 & 0.497 & 25 & 0.567$^\dagger$ & 0.497 \\
\midrule
IIIT-Delhi Academic & 2015 & \textbf{0.393} & [0.28, 0.50] & 3 & 0.477 & 25 & 0.714$^\dagger$ & 0.616$^\dagger$ \\
IIIT-Delhi Academic & 2016 & \textbf{0.636} & [0.50, 0.75] & 5 & 0.580 & 25 & 0.732 & 0.631 \\
IIIT-Delhi Academic & 2017 & \textbf{0.422} & [0.30, 0.54] & 5 & 0.595$^\dagger$ & Gauss. & 0.801$^\dagger$ & 0.648$^\dagger$ \\
IIIT-Delhi Boys & 2014 & \textbf{0.501} & [0.38, 0.61] & 4 & 0.554 & 25 & 0.746$^\dagger$ & 0.642$^\dagger$ \\
IIIT-Delhi Campus & 2017 & \textbf{0.432} & [0.29, 0.58] & 4 & 0.482 & 25 & 0.636$^\dagger$ & 0.553 \\
IIIT-Delhi Dining & 2017 & \textbf{0.714} & [0.61, 0.81] & 15 & 0.744 & 25 & 0.774 & 0.665 \\
IIIT-Delhi Facilities & 2015 & \textbf{0.639} & [0.55, 0.71] & 7 & 0.705 & Gauss. & 0.875$^\dagger$ & 0.705 \\
IIIT-Delhi Facilities & 2016 & \textbf{0.621} & [0.52, 0.72] & 5 & 0.633 & Gauss. & 0.856$^\dagger$ & 0.690 \\
IIIT-Delhi Facilities & 2017 & \textbf{0.739} & [0.62, 0.85] & 10 & 0.753 & Gauss. & 0.876$^\dagger$ & 0.706 \\
IIIT-Delhi Girls & 2014 & \textbf{0.693} & [0.59, 0.78] & 4 & 0.589$^\dagger$ & Gauss. & 0.848$^\dagger$ & 0.684 \\
IIIT-Delhi Girls & 2017 & \textbf{0.818} & [0.74, 0.89] & 7 & 0.739 & Gauss. & 0.914$^\dagger$ & 0.739 \\
IIIT-Delhi Library & 2014 & \textbf{0.480} & [0.36, 0.59] & 4 & 0.567 & Gauss. & 0.816$^\dagger$ & 0.659$^\dagger$ \\
IIIT-Delhi Library & 2016 & \textbf{0.509} & [0.41, 0.61] & 4 & 0.504 & 25 & 0.670$^\dagger$ & 0.579 \\
\bottomrule
\end{tabular}
}
\caption{What selecting the copula's degrees of freedom contributes. Realised horizon exceedance at $H=64$ on the held-out month with its day-block bootstrap interval, and the copula's prediction under three choices of $\nu$: matched to horizon exceedance on the validation block (the paper's choice; columns 5--6), maximising the $t$-copula pseudo-likelihood on the same block (columns 7--8), and held at the Phoenix value $\nu=7$ with no refitting (last column). $\dagger$ marks a prediction outside the realised interval. The tail-matched choice lies inside on 12 of 14 windows (mean absolute error $0.065$). The likelihood criterion selects $\nu \geq 25$ or the Gaussian on every window, over-predicts on 14 of 14 and lies inside on 2: it is dominated by the body of the distribution, where the families agree, and cannot see the upper-tail dependence the horizon event is made of. The Phoenix value applied unchanged to the 13 Indian windows lies inside on 9 of 13 (mean absolute error $0.098$); the per-window selection buys accuracy, not the agreement itself. The independence figure lies inside on none.}
\label{tab:copula-df}
\end{table}


Two further caveats belong here rather than in a footnote, and one that was a
caveat has been measured. The copula's degrees of freedom are selected on the
validation block by matching predicted to realised horizon exceedance, which is
the quantity the copula is later judged on; the split is honest, but the
criterion knows what it will be graded on. Table~\ref{tab:copula-df} bounds what
that choice contributes. Holding $\nu$ at the Phoenix value of $7$ and applying
it unchanged to the thirteen Indian windows, the copula still lies inside the
realised interval on nine of thirteen, with mean absolute error $0.098$ against
$0.063$ when $\nu$ is selected per window: the selection buys accuracy, not the
agreement. The textbook alternative --- choosing $\nu$ by the $t$-copula
pseudo-likelihood on the same validation block --- selects $\nu \geq 25$ or the
Gaussian on every one of the fourteen windows, over-predicts on all fourteen,
and lies inside the interval on two. The likelihood is dominated by the body of
the distribution, where the two families agree; the horizon event is a statement
about the joint upper tail, where they do not. That is why the one free
parameter is fitted to the decision-relevant statistic rather than to the
likelihood, and it is a warning for anyone fitting a dependence model for a
tail event by the default criterion. Base-load and PV scenarios are drawn
independently, while cloud physically raises cooling load at the same moment it
cuts PV output, so the true joint law has worse afternoons than these scenarios
contain. And the horizon is sixteen hours while the demand charge bills a
monthly maximum over roughly $2{,}880$ intervals: a joint guarantee over the
window is strictly stronger than a marginal one and strictly weaker than a
monthly one. Extending it needs the copula fitted over a month-length horizon,
which is a $2{,}880$-square correlation matrix and needs structure rather than a
raw estimate.

The honest summary is not that the guarantee is restored exactly. It is that a
risk level which was previously an \emph{emergent property of horizon length}
--- unstated, unsettable, and equal to $0.106$ at the tight target on the
office and $0.297$ on the campus feed under the incumbent marginal controller
--- is now a parameter with a monotone response on both, conservative on the
one whose marginals are calibrated, at a cost of \rupee{3{,}339} on the month
there and \rupee{3{,}329} \emph{saved} on the campus, where the tight target
binds so rarely that the scenario controller buys its risk level out of
headroom the marginal one was wasting.

\section{What degrades with aggregation, and what does not}
\label{sec:calibration}

One protocol --- 15 months train, 2 months validate, 1 month test, temporal
throughout, conformal calibration on the validation block --- run unchanged over
thirty supplies. Tables~\ref{tab:climate}--\ref{tab:china} give the five arms.

\begin{table}[t]
\centering
\footnotesize
\setlength{\tabcolsep}{3.5pt}
\resizebox{\linewidth}{!}{%
\begin{tabular}{llrrrrrr}
\toprule
Site & Country & CDD sh. & HVAC sh. & Seas.\ naive & Ours & Skill & Cov 90\% \\
\midrule
Phoenix, United States & United States & 0.87 & 0.38 & 4.529 & 3.484 & \textbf{+0.231} & 0.884 \\
Austin, United States & United States & 0.75 & 0.13 & 1.352 & 1.061 & \textbf{+0.215} & 0.906 \\
Washington DC, United States & United States & 0.31 & 0.23 & 4.698 & 2.964 & \textbf{+0.369} & 0.900 \\
Minneapolis, United States & United States & 0.11 & 0.13 & 21.875 & 10.792 & \textbf{+0.507} & 0.916 \\
Berkeley, United States & United States & 0.05 & 0.26 & 5.855 & 3.033 & \textbf{+0.482} & 0.873 \\
London, United Kingdom & United Kingdom & 0.05 & 0.04 & 4.106 & 2.061 & \textbf{+0.498} & 0.910 \\
Dublin, Ireland & Ireland & 0.00 & 0.04 & 7.636 & 3.451 & \textbf{+0.548} & 0.918 \\
\bottomrule
\end{tabular}
}
\caption{Climate arm: demographic held fixed at Education, climate and country varied across seven sites and three countries.}
\label{tab:climate}
\end{table}

\begin{table}[t]
\centering
\footnotesize
\setlength{\tabcolsep}{3.5pt}
\resizebox{\linewidth}{!}{%
\begin{tabular}{llrrrrrr}
\toprule
Site & Country & CDD sh. & HVAC sh. & Seas.\ naive & Ours & Skill & Cov 90\% \\
\midrule
Phoenix, United States & United States & 0.87 & 0.24 & 6.213 & 2.419 & \textbf{+0.611} & 0.888 \\
Washington DC, United States & United States & 0.31 & 0.05 & 49.625 & 34.038 & \textbf{+0.314} & 0.824 \\
Minneapolis, United States & United States & 0.11 & 0.24 & 4.784 & 3.009 & \textbf{+0.371} & 0.914 \\
London, United Kingdom & United Kingdom & 0.05 & 0.01 & 13.893 & 20.322 & \textbf{-0.463} & 0.930 \\
Dublin, Ireland & Ireland & 0.00 & 0.03 & 4.150 & 2.020 & \textbf{+0.513} & 0.912 \\
\bottomrule
\end{tabular}
}
\caption{Office arm: the same ladder in a second usage class, run as a control on the climate arm. The London office is the study's worst row and is reported unchanged.}
\label{tab:office}
\end{table}

\begin{table}[t]
\centering
\footnotesize
\setlength{\tabcolsep}{3.5pt}
\resizebox{\linewidth}{!}{%
\begin{tabular}{llrrrrrr}
\toprule
Usage class & Country & CDD sh. & HVAC sh. & Seas.\ naive & Ours & Skill & Cov 90\% \\
\midrule
Entertainment/public assembly & United States & 0.31 & 0.31 & 4.193 & 2.038 & \textbf{+0.514} & 0.919 \\
Healthcare & United States & 0.31 & 0.21 & 2.907 & 1.702 & \textbf{+0.415} & 0.902 \\
Retail & United States & 0.31 & 0.25 & 7.681 & 4.734 & \textbf{+0.384} & 0.897 \\
Warehouse/storage & United States & 0.31 & 0.06 & 1.635 & 1.189 & \textbf{+0.273} & 0.942 \\
Lodging/residential & United States & 0.31 & 0.28 & 2.726 & 2.102 & \textbf{+0.229} & 0.901 \\
Public services & United States & 0.31 & 0.03 & 3.882 & 3.680 & \textbf{+0.052} & 0.887 \\
\bottomrule
\end{tabular}
}
\caption{Demographic arm: climate held fixed at Washington DC, demographic varied across the eight usage classes that pass the selection rule at that site. Six rows appear here; Education and Office were also run in this arm and are shown in Tables~\ref{tab:climate} and~\ref{tab:office} instead, since a building is one supply however many arms it appears in.}
\label{tab:demographic}
\end{table}

\begin{table}[t]
\centering
\footnotesize
\setlength{\tabcolsep}{3.5pt}
\resizebox{\linewidth}{!}{%
\begin{tabular}{llrrrrrr}
\toprule
Series & Country & CDD sh. & HVAC sh. & Seas.\ naive & Ours & Skill & Cov 90\% \\
\midrule
India (Delhi) & India & 0.89 & -- & 72.499 & 65.437 & \textbf{+0.097} & 0.762 \\
Spain & Spain & 0.28 & -- & 594.320 & 226.544 & \textbf{+0.619} & 0.813 \\
France & France & 0.08 & -- & 942.571 & 342.809 & \textbf{+0.636} & 0.829 \\
Germany & Germany & 0.06 & -- & 1344.698 & 386.465 & \textbf{+0.713} & 0.837 \\
United Kingdom & United Kingdom & 0.02 & -- & 451.549 & 378.932 & \textbf{+0.161} & 0.822 \\
Ireland & Ireland & 0.00 & -- & 29.812 & 16.624 & \textbf{+0.442} & 0.911 \\
\bottomrule
\end{tabular}
}
\caption{National arm, metered: system-level demand in six countries. Delhi is 15-minute native resolution from the Delhi State Load Despatch Centre; the European series are ENTSO-E hourly via Open Power System Data.}
\label{tab:national}
\end{table}

\begin{table}[t]
\centering
\footnotesize
\setlength{\tabcolsep}{3.5pt}
\resizebox{\linewidth}{!}{%
\begin{tabular}{llrrrrrrr}
\toprule
Series & Country & Load--temp. & CDD sh. & HVAC sh. & Seas.\ naive & Ours & Skill & Cov 90\% \\
\midrule
Hainan (Haikou) & China & +0.10 & 0.97 & -- & 131.650 & 93.878 & \textbf{+0.287} & 0.960 \\
Guangdong (Guangzhou) & China & +0.37 & 0.85 & -- & 2357.241 & 3097.316 & \textbf{-0.314} & 0.919 \\
Shanghai & China & +0.07 & 0.44 & -- & 445.551 & 352.338 & \textbf{+0.209} & 0.894 \\
Beijing & China & -0.28 & 0.29 & -- & 451.459 & 387.453 & \textbf{+0.142} & 0.905 \\
Yunnan (Kunming) & China & -0.35 & 0.14 & -- & 251.362 & 197.905 & \textbf{+0.213} & 0.703 \\
Heilongjiang (Harbin) & China & -0.47 & 0.12 & -- & 148.038 & 191.583 & \textbf{-0.294} & 0.565 \\
\bottomrule
\end{tabular}
}
\caption{The China arm, added after the study was complete on a constraint imposed from outside it: six provinces spanning a wider climate range than the entire European panel, in a developing economy, at one calendar year and therefore a three-month training block rather than fifteen. Mean skill is $+0.040$ with 2 negative rows. \textbf{That column should not be read as a statement about the method}, for the reason established in Table~\ref{tab:china-audit}: the series is constructed, not metered, and its construction pins our model and the baseline it is scored against to the same two numbers per day. Coverage does not run through the baseline and is readable; skill is not.}
\label{tab:china}
\end{table}


\begin{table}[t]
\centering
\small
\resizebox{\linewidth}{!}{%
\begin{tabular}{lrrrr}
\toprule
Population & $n$ & mean cov. & worst & $<0.85$ \\
\midrule
Buildings (BDG2), metered & 18 & \textbf{0.901} & 0.824 & 1 \\
System demand, metered & 6 & \textbf{0.829} & 0.762 & 5 \\
System demand, reconstructed & 6 & \textbf{0.824} & 0.565 & 2 \\
\midrule
\emph{India, one city (I-BLEND + Delhi SLDC), windows} &  &  &  &  \\
\quad buildings, native 15-min & 12 & \textbf{0.910} & 0.897 & 0 \\
\quad campus feed (one HT consumer) & 1 & \textbf{0.873} & 0.873 & 0 \\
\quad city (Delhi SLDC) & 1 & \textbf{0.762} & 0.762 & 1 \\
\bottomrule
\end{tabular}
}
\caption{Empirical coverage of the nominal 90\% interval after conformal calibration, \emph{as the benchmark calibrated it}: split conformal plus an adaptive layer at the step $\gamma=0.35$ in the units of each series. Read this way the guarantee very nearly holds on individual buildings and fails systematically on system-level demand. The third row is the replication: six Chinese provinces, a different country, a different year and a different data-generating process, reproduce the system-level failure to within 0.005 of the metered panel and contain the worst row in the study (Heilongjiang, 0.565). The lower panel is the same effect inside one city: building, campus and city demand in Delhi, one weather feed, native 15-minute resolution throughout, and coverage falls monotonically with each step up the ladder. Table~\ref{tab:aci-step} re-reads every row of this table with the adaptive step stated in the interval's units, and the ladder does not survive it.}
\label{tab:calibration}
\end{table}

\begin{table}[t]
\centering
\footnotesize
\setlength{\tabcolsep}{3.5pt}
\resizebox{\linewidth}{!}{%
\begin{tabular}{lrrrrrrr}
\toprule
Population & $n$ & $\kappa$ of $\gamma{=}0.35$ & split & ACI $\gamma{=}0.35$ & ACI $\kappa{=}0.006$ & worst & $<0.85$ \\
\midrule
Buildings (BDG2), metered & 18 & 0.0113 & 0.820 & 0.901 & \textbf{0.898} & 0.862 & 0 \\
System demand, metered & 6 & 0.0001 & 0.791 & 0.829 & \textbf{0.914} & 0.902 & 0 \\
System demand, reconstructed & 6 & 0.0001 & 0.785 & 0.824 & \textbf{0.925} & 0.902 & 0 \\
\midrule
\emph{India, one city, windows} &  &  &  &  &  &  &  \\
\quad buildings, native 15-min & 12 & 0.0338 & 0.868 & 0.910 & \textbf{0.908} & 0.874 & 0 \\
\quad campus feed (one HT consumer) & 1 & 0.0016 & 0.819 & 0.873 & \textbf{0.888} & 0.888 & 0 \\
\quad city (Delhi SLDC) & 1 & 0.0006 & 0.605 & 0.762 & \textbf{0.904} & 0.904 & 0 \\
\bottomrule
\end{tabular}
}
\caption{The ladder of Table~\ref{tab:calibration} re-read without the confound. The paper's forecaster refitted once per supply and the same predictions calibrated three ways: split conformal alone; split plus ACI at the absolute step $\gamma=0.35$ every benchmark row used; and split plus ACI at the building audit's step as a fraction of the interval width, $\kappa=0.006$, the width taken from the validation block so that nothing from the test month is used. The third column is what the absolute step amounted to on each population, as a fraction of its interval (median): a working adaptive layer on buildings and a nearly inert one on system demand. Mean coverage of the nominal $0.90$ interval on the test month, then the worst row and the count below $0.85$ under the relative step. The building--system gap is $+0.029$ under split conformal alone, $+0.073$ under the absolute step the benchmark used, and $-0.016$ under the step stated in the interval's units.}
\label{tab:aci-step}
\end{table}


\paperfig{fig_calibration.pdf}{One dot per supply, twice. Top, as the benchmark calibrated it: the
nominal $90\%$ interval very nearly holds on individual buildings and fails
once demand is aggregated to system level, on the metered panel and the
reconstructed one alike, two populations sharing no country, no year and no
data-generating process whose means differ by $0.005$. Bottom, the same
supplies with the adaptive layer's step stated as a fraction of the interval
width rather than as a number in the series' units: every tier is at the level
and the ladder is gone.}{fig:calibration}

Table~\ref{tab:calibration} and the upper panel of Figure~\ref{fig:calibration}
are the benchmark's own reading, and they were the third contribution of the
first draft of this paper. A nominal $90\%$ interval holds at building level
(mean $0.901$; one row of eighteen below $0.85$) and fails at system level
(mean $0.829$ metered, $0.824$ reconstructed; five of six and two of six below
$0.85$ respectively); Delhi is the worst metered row at $0.762$ and
Heilongjiang the worst row anywhere at $0.565$. The audit of the adaptive layer
at system scale, later in this section, is where that reading came apart.
Every row of the benchmark was calibrated with split conformal and an adaptive
layer whose step, $\gamma=0.35$, is a number in the units of the series: on a
building whose intervals are tens of kilowatts wide it is a working layer, and
on a grid whose intervals are thousands of megawatts wide it is one that does
not move. The building rows carried a remedy the system rows did not, and the
ladder measured the remedy's absence.

Table~\ref{tab:aci-step} is the same measurement without the confound: the
paper's forecaster refitted once on every supply and the same predictions
calibrated three ways. Under split conformal alone the tiers sit within $0.03$
of each other --- $0.820$ on buildings, $0.791$ on metered system demand,
$0.785$ reconstructed --- and eleven of the eighteen buildings are below
$0.85$. The exchangeability failure is universal, not a property of
aggregation. Under the adaptive layer with its step stated as a fraction of the
interval width --- $\kappa=0.006$, the width read from the validation block so
that nothing from the test month is used --- every tier is at the level,
$0.898$, $0.914$ and $0.925$, and no supply in the study is below $0.85$:
Heilongjiang is $0.937$ and Delhi $0.904$. The building--system gap is
$+0.029$, $+0.073$ and $-0.016$ under the three calibrations. What degraded
with aggregation was the step.

This is consequential rather than curious, because the controller's $q95$ is
read from the same object, and the consequence has changed sign since the first
draft. The hazard is not a tier the guarantee cannot reach. It is a parameter
that goes quietly inert when the layer is moved to a larger meter, while the
coverage it reports looks like a property of the meter.

\paragraph{The mechanism.} Split conformal guarantees coverage when the
calibration and test blocks are exchangeable
\citep{vovk2005algorithmic,romano2019cqr}. A Delhi May is not exchangeable with
a Delhi June, because the monsoon arrives between them and changes the
relationship between temperature and load. This is an assumption failure, not a
tuning failure, and it is diagnosable rather than speculative: our own audit
partitions May into six disjoint calibration blocks and finds coverage varying
from $0.710$ to $0.861$ with a block-bootstrap standard error of $0.031$,
failing every acceptance criterion stated in advance.\footnote{The bootstrap
resamples thirty whole test days rather than $184{,}320$ overlapping 15-minute
forecasts, because those are not $184{,}320$ independent observations. The naive
standard error is under $0.001$ and would pass any model ever built.}

\paperfig{fig_aci.pdf}{The remedy, measured at both tiers. Twelve monthly folds, each
trained strictly on its own past and calibrated on the thirty days immediately
before it. Top, the Phoenix office: raw quantiles are not calibrated out of
sample at all; split conformal --- the layer with the theorem attached ---
narrows the failure and does not remove it, bottoming out at $0.447$ in a month
whose ceiling was being defended against a bound reality broke through a third
of the time; adaptive conformal inference holds $0.867$--$0.928$ in every month.
Bottom, the Delhi city feed under the identical protocol moved five years back
onto its archive. Red is ACI at the building's step $\gamma=0.35$, a number in
the series' units; blue is the same step carried over as the building's
fraction of the interval width, $\kappa=0.006$, with nothing tuned on Delhi.}{fig:aci}

\begin{table}[t]
\centering
\footnotesize
\setlength{\tabcolsep}{3.5pt}
\resizebox{\linewidth}{!}{%
\begin{tabular}{lrrrrrrrr}
\toprule
Series & $\gamma$ & $\kappa$ & split, by month & in band & ACI, by month & in band & split & ACI \\
\midrule
Phoenix office, building, 2016--17 & 0.35 & 0.0060 & 0.447--0.936 & 52\% & 0.867--0.928 & 87\% & 0.770 & 0.885 \\
Delhi city, system, 2011--12 & 0.35 & 0.0005 & 0.628--0.980 & 48\% & 0.765--0.969 & 69\% & 0.358 & 0.733 \\
Delhi city, system, 2011--12 & 4.43 & 0.0060 & 0.628--0.980 & 48\% & 0.874--0.935 & 100\% & 0.358 & 0.887 \\
\bottomrule
\end{tabular}
}
\caption{The adaptive layer at each tier it has been run at. Walk-forward year of twelve monthly folds, each trained strictly on the past and calibrated on the thirty days before it; coverage of the nominal $0.90$ interval by month (range) and the share of the rolling 30-day curve inside $0.85$--$0.95$. The last two columns are coverage after a level and volatility shift injected into a model frozen six months earlier, which static calibration cannot respond to by construction. $\gamma$ is the ACI step in the units of the series and $\kappa$ the same step as a fraction of the mean split-conformal interval width at the audited lead. At the building, split conformal spends 52\% of the year in band and ACI 87\%. Carried to the city as the same absolute $\gamma$, the step is 13$\times$ smaller relative to the interval and ACI reaches 69\% against split's 48\%. Carried as the same $\kappa$, with nothing tuned on Delhi, ACI spends 100\% of the year in band, by-month coverage 0.874--0.935, and settles at 0.887 after the shift against split's 0.358.}
\label{tab:aci}
\end{table}


\paragraph{The remedy, at building scale.} Table~\ref{tab:aci} and
Figure~\ref{fig:aci} run a walk-forward year. Split conformal ranges over
$0.447$--$0.936$ by month and spends $52\%$ of the year inside a
$0.85$--$0.95$ band; adaptive conformal inference \citep{gibbs2021adaptive}
ranges over $0.867$--$0.928$ and spends $87\%$. Under a frozen model and a
synthetic distribution shift the same ordering holds: split conformal settles at
$0.770$ coverage and never recovers within the run, ACI at $0.885$. The layer
with the finite-sample theorem is not the layer that delivers the level here,
and that is worth stating plainly because it inverts the usual presentation
order.

\begin{table}[t]
\centering
\small
\resizebox{\linewidth}{!}{%
\begin{tabular}{lrrrrrr}
\toprule
$\kappa$ & in band & by month & width & in band & by month & width \\
\midrule
0 (split only) & 52.2\% & 0.447--0.936 & 58 & 48.1\% & 0.628--0.980 & 739 \\
0.001 & 61.4\% & 0.785--1.000 & 67 & 72.9\% & 0.825--0.960 & 756 \\
0.002 & 81.8\% & 0.816--0.992 & 64 & 89.7\% & 0.840--0.956 & 735 \\
0.005 & 85.2\% & 0.860--0.941 & 61 & 99.6\% & 0.869--0.939 & 698 \\
0.01 & 93.1\% & 0.878--0.932 & 59 & 100.0\% & 0.882--0.921 & 672 \\
0.02 & 100.0\% & 0.890--0.924 & 53 & 100.0\% & 0.894--0.908 & 629 \\
0.05 & 100.0\% & 0.896--0.909 & 46 & 100.0\% & 0.897--0.904 & 560 \\
0.1 & 100.0\% & 0.897--0.906 & 41 & 100.0\% & 0.898--0.902 & 495 \\
\bottomrule
\end{tabular}
}
\caption{ACI replayed over the saved walk-forward year at a grid of step sizes, stated as a fraction $\kappa$ of the mean split-conformal interval width at the audited lead (58, 739 in the series' units for the Phoenix office and Delhi city respectively; left block Phoenix, right block Delhi). Nothing is retrained: ACI is a pass over the saved split-conformal predictions in time order. In band is the share of the rolling 30-day coverage curve inside $0.85$--$0.95$; by month is the range of monthly coverage; width is the mean adaptive interval. The audit's absolute default $\gamma=0.35$ is $\kappa = 0.0060, 0.0005$ on the two series. Above $\kappa \approx 0.01$ the interval is a fast tracker rather than a forecast interval: the recursion pins the long-run rate whatever the model, and the width it reports is no longer the forecaster's.}
\label{tab:aci-gamma}
\end{table}


\paragraph{The remedy, at system scale.} The same audit on the Delhi city feed
--- the tier where the failure is systematic, and the tier India sits in --- is
the second and third rows of Table~\ref{tab:aci} and the lower panel of
Figure~\ref{fig:aci}, the 2016--17 protocol moved five years back onto the
2011--12 archive so that the calendar position is identical. Split conformal
fails harder here than at the building. Six disjoint May calibration blocks
give June coverage from $0.344$ to $0.728$, mean $0.524$ with a day-block
bootstrap standard error of $0.044$; the walk-forward year ranges
$0.628$--$0.980$ by month and spends $48\%$ in band; after the frozen shift it
settles at $0.358$. The adaptive layer, carried over with the step it had at
the building, $\gamma=0.35$, reaches $69\%$ in band and $0.733$ after the
shift: better, and not the level.

But $\gamma$ is a number in the units of the series. On a building whose
intervals are $58$\,kW wide it moves the bound by $0.6\%$ of the interval per
step; on a city whose intervals are $739$\,MW wide, by $0.05\%$. The same
setting is a layer twelve times slower, and Table~\ref{tab:aci-gamma} shows
that this is the whole of the difference: ACI replayed over the saved year at a
grid of steps stated as a fraction $\kappa$ of the interval width, nothing
retrained, reaches the level at the city as it does at the building. Carried
over as the building's $\kappa=0.006$ rather than its $\gamma$ --- nothing
tuned on Delhi --- the adaptive layer spends $100\%$ of the year in band with
by-month coverage $0.874$--$0.935$, holds the monsoon June that the benchmark
row reports at $0.762$ at $0.895$, and settles at $0.887$ after the shift,
where the building settled at $0.885$. The level is restored at system scale.
What did not transfer was the parameter's units, and a practitioner moving the
layer from a building to a feeder would hit exactly this. One caveat is due:
above $\kappa\approx0.01$ the interval is a fast tracker rather than a
forecast interval, since the recursion pins the long-run exceedance rate
whatever the model and the width it reports is no longer the forecaster's; the
sweep is read for the transfer, not for the numbers at its top.

\paragraph{Selection is regional, and we found that out the hard way.} The
rule's final clause --- positive load--temperature correlation --- is a sanity
check in Phoenix, where every candidate is cooling-dominated, and a
site-destroying filter in Minneapolis or Dublin, where buildings consume
electricity when it is cold. Seventy-seven buildings across nine sites pass every
clause but that one. Separately, the two hottest sites in the corpus, the closest
available analogues to an Indian metro climate, yield no admissible building at
all for want of coverage in the split window. Both facts were discovered by
applying the rule, not anticipated when writing it.

\section{Skill and flexibility are independent axes}
\label{sec:axes}

\begin{table}[t]
\centering
\small
\resizebox{\linewidth}{!}{%
\begin{tabular}{lrrl}
\toprule
Correlation against forecast skill & $r$ & $n$ & leave-one-out range \\
\midrule
\emph{Panel A --- 18 BDG2 buildings, three countries} &  &  &  \\
\quad cooling-degree-day share & +0.010 & 18 & -0.308 to +0.077 \\
\quad controllable (HVAC) fraction & +0.283 & 18 & +0.029 to +0.378 \\
\quad mean outdoor temperature & -0.012 & 18 & -0.266 to +0.051 \\
\quad median load & -0.014 & 18 & -0.189 to +0.120 \\
\midrule
\emph{Panel B --- 6 Chinese provinces, independent replication} &  &  &  \\
\quad cooling-degree-day share & +0.054 & 6 & -0.377 to +0.582 \\
\quad mean outdoor temperature & +0.216 & 6 & -0.290 to +0.782 \\
\quad load--temperature correlation & -0.060 & 6 & -0.619 to +0.713 \\
\midrule
\emph{controllable fraction} vs climate (not vs skill) & \textbf{+0.489} & 18 & -- \\
\bottomrule
\end{tabular}
}
\caption{Forecast skill is uncorrelated with climate, controllable fraction, temperature and load size, and every leave-one-out range straddles or nearly straddles zero --- which at these sample sizes is what a null looks like. Panel B is the replication on a panel the null was not fitted to. The final row is the one relationship in the study that is real, and note what it is between: available flexibility tracks climate strongly, while predictability does not track anything. Skill and flexibility are independent deployment axes.}
\label{tab:correlations}
\end{table}


\paperfig{fig_null.pdf}{(a) The null. (b) The same null on a panel it was not fitted to, in a
different country, year and data-generating process; panels (a) and (b) share a
vertical scale. (c) The same eighteen buildings, the same horizontal axis, and
the one relationship in the study that is real. How much load there is to shift
tracks climate. How well it can be predicted does not.}{fig:null}

We expected forecast skill --- measured with pinball loss, a strictly proper
scoring rule \citep{gneiting2007strictly}, against each series' own
seasonal-naive baseline --- to increase with cooling load, on the reasoning that
weather-driven demand is the learnable part. Table~\ref{tab:correlations} and
Figure~\ref{fig:null} report the opposite: a null. Every correlation against
skill is indistinguishable from zero and every leave-one-out range straddles or
nearly straddles the sign line, which at these sample sizes is what a null
result looks like. It replicates at $r=+0.054$ on the Chinese provincial panel,
which was added afterwards and spans a wider climate range than the whole
European set.

The controllable fraction, by contrast, tracks climate strongly ($r=+0.489$),
running from $37.6\%$ of the meter in Phoenix to $0.8\%$ in London.

The deployment consequence is a screening rule. Do not ask whether a site's load
is predictable; ask how much of it is weather-driven. A Dublin office is highly
predictable (skill $+0.513$) with roughly $3\%$ controllable load, and there is
nothing to sell there. Conflating the axes inverts the decision. For the target
market this is favourable: Delhi demand correlates $+0.59$ with outdoor
temperature against $-0.64$ for France, placing India at the useful end of the
axis that matters.

\section{The new country, and what it turned out to measure}
\label{sec:china}

The study was complete when a constraint was imposed from outside it: add a
country. China was chosen because six provinces spanning Hainan to Heilongjiang
cover a wider climate range than the entire European panel and do so in a
developing economy, which is the comparison India actually needs. The protocol
adaptation it forced is stated in Table~\ref{tab:china} and is a single change:
the source is one calendar year, so the fifteen-month training block does not
exist and the calendar \emph{position} of the split is preserved instead ---
validate on April and May, test on June, as in every other row --- leaving three
months of training rather than fifteen. That biases against us, which is the safe
direction to err.

\begin{table}[t]
\centering
\small
\resizebox{\linewidth}{!}{%
\begin{tabular}{lrrlrrr}
\toprule
Province & Days & Shapes & Days/shape & Numbers & Values & Max err.\ \% \\
\midrule
Hainan (Haikou) & 365 & \textbf{2} & 115/250 & 778 & 8,760 & 0.000000 \\
Guangdong (Guangzhou) & 365 & \textbf{2} & 115/250 & 778 & 8,760 & 0.000000 \\
Shanghai & 365 & \textbf{2} & 115/250 & 778 & 8,760 & 0.000000 \\
Beijing & 365 & \textbf{2} & 115/250 & 778 & 8,760 & 0.000000 \\
Yunnan (Kunming) & 365 & \textbf{2} & 115/250 & 778 & 8,760 & 0.000000 \\
Heilongjiang (Harbin) & 365 & \textbf{2} & 115/250 & 778 & 8,760 & 0.000000 \\
\bottomrule
\end{tabular}
}
\caption{Provenance audit of the China arm, run against the data rather than taken from its documentation. Every calendar day's 24-hour profile is min--max normalised and the distinct shapes counted. \emph{Shapes} is how many survive; \emph{Numbers} is how many values are needed to rebuild the year from them; \emph{Values} is how many the year contains. A whole year of hourly provincial demand is 778 numbers --- 2 normalised day-shapes and one (min, max) pair per day --- reproducing all 8,760 values to within 2e-09 of their magnitude, which is the precision the source is published at rather than an approximation we introduced. Compression 11.26:1. The published method is exactly what the authors state it is. The consequence for this study is that forecast skill, a ratio against a seasonal-naive baseline, is uninformative on these series: the shape the baseline has to guess is constant, so both models reduce to predicting the same two administrative numbers per day.}
\label{tab:china-audit}
\end{table}


\paragraph{The audit.} This is the only tier-2 source whose documentation
describes a construction rather than a measurement. Wu and Kan take each day's
maximum and minimum from National Development and Reform Commission reporting,
trace a representative workday profile and a representative holiday profile off
published load curves with WebPlotDigitizer, and rescale one of those two
profiles to each day's envelope \citep{wu2023china}. Taking that on trust is
exactly what this project does not do elsewhere, so we checked it against the
data: min--max normalise every calendar day's 24-hour profile and count the
distinct shapes in a year.

The answer is two, in every province. Table~\ref{tab:china-audit} reports it:
$778$ numbers --- two normalised day-shapes and $365$ (min, max) pairs ---
reproduce all $8{,}760$ hourly values to within $2\times10^{-9}$ of their
magnitude, which is the precision the source is published at rather than an
approximation we introduced, at a compression of better than eleven to one. The
documentation is accurate.

\paragraph{What that disqualifies.} Skill in this study is a ratio against each
series' own seasonal-naive baseline. When every workday shares one shape,
seasonal naive recovers that shape exactly and its entire error is the daily
envelope --- and so is ours. Both models collapse onto predicting the same two
administrative numbers per day, the ratio between them collapses toward zero,
and the observed mean skill of $+0.040$ with two negative rows is what that
predicts rather than a finding about a forecaster. The China skill column is
reported for completeness and is not evidence either way about the method.

This costs us the arm's most quotable result. Guangdong is the closest structural
analogue to Delhi anywhere in the study --- cooling share $0.85$ against Delhi's
$0.89$, load--temperature correlation $+0.37$ against $+0.59$, a comparably peaky
profile --- and it returns $-0.314$, the second-worst skill in the benchmark.
Read naively that corroborates the Indian result on independent 2018 data from
another country and retires the objection that Delhi's weakness is an artefact of
a decade-old archive. It does not, because the number is a property of the
construction. We report the temptation as well as the retraction, because the
inference is one a reader would otherwise make on our behalf.

\paragraph{What survives.} Coverage is a property of interval width against
realised error and does not run through the baseline, so the tier-2 result can be
read on this arm --- and it replicates twice over. Under the benchmark's
calibration the six provinces average $0.824$ against $0.829$ on the metered
panel, on an error process unlike any other in the study, with the worst row
anywhere at $0.565$; under the adaptive step stated in the interval's units
they average $0.925$ against $0.914$, and the worst row, Heilongjiang, is
$0.937$ (Table~\ref{tab:aci-step}). Two populations sharing no country, no
year and no data-generating process landing within $0.005$ of each other under
the inert layer and within $0.011$ under the working one is stronger evidence
for both readings than a larger metered sample would have been: the failure is
universal, and so is the remedy. The null of Section~\ref{sec:axes} replicates
too. The wider coverage spread on this panel under the benchmark's calibration
($0.565$--$0.960$ against $0.762$--$0.911$) is consistent with the three-month
calibration history and is not separately identified.

\paragraph{What the exercise is worth.} A constraint imposed late produced no new
capability and one new measurement, which is a characterisation of the dataset
rather than of the method, and which removed a result we would have liked to
keep. That is the outcome we would want from an audit, and it is reported at the
same size as everything else.

\section{India in comparison}
\label{sec:india}

\begin{table}[t]
\centering
\footnotesize
\setlength{\tabcolsep}{3.5pt}
\resizebox{\linewidth}{!}{%
\begin{tabular}{llrrrrrr}
\toprule
Series & Provenance & Native min & Cooling share & Load--temp. corr. & $p99/$med & Skill & Cov 90\% \\
\midrule
Hainan (Haikou) & recon. & 60 & 0.97 & +0.10 & 1.32 & +0.287 & 0.960 \\
\textbf{India (Delhi)} & metered & 15 & 0.89 & +0.59 & 1.61 & +0.097 & 0.762 \\
Guangdong (Guangzhou) & recon. & 60 & 0.85 & +0.37 & 1.47 & -0.314 & 0.919 \\
Shanghai & recon. & 60 & 0.44 & +0.07 & 1.64 & +0.209 & 0.894 \\
Beijing & recon. & 60 & 0.29 & -0.28 & 1.59 & +0.142 & 0.905 \\
Spain & metered & 60 & 0.28 & +0.17 & 1.31 & +0.619 & 0.813 \\
Yunnan (Kunming) & recon. & 60 & 0.14 & -0.35 & 1.24 & +0.213 & 0.703 \\
Heilongjiang (Harbin) & recon. & 60 & 0.12 & -0.47 & 1.31 & -0.294 & 0.565 \\
France & metered & 60 & 0.08 & -0.64 & 1.58 & +0.636 & 0.829 \\
Germany & metered & 60 & 0.06 & -0.10 & 1.32 & +0.713 & 0.837 \\
United Kingdom & metered & 60 & 0.02 & -0.13 & 1.40 & +0.161 & 0.822 \\
Ireland & metered & 60 & 0.00 & -0.07 & 1.37 & +0.442 & 0.911 \\
\bottomrule
\end{tabular}
}
\caption{India against every other system-level series in the study, ordered by cooling-degree-day share. Among metered supplies Delhi is the extreme of both halves of the verdict: the highest cooling share, the strongest positive load--temperature relationship, the peakiest profile, and simultaneously the lowest forecast skill and the worst interval coverage. The six Chinese provinces, added after the study was written, displace India from three of those four extremes on the face of the table --- and the provenance column is why that displacement is reported rather than acted on. Guangdong is the closest structural analogue to Delhi in the study and returns the second-worst skill in it, which reads like corroboration and is not: Section~\ref{sec:china} shows the Chinese skill column is an artefact of how the data was built. Delhi remains the only series at native 15-minute resolution, the cadence the controller runs at and the cadence Indian demand charges are assessed on.}
\label{tab:india}
\end{table}


Table~\ref{tab:india} is the comparison the benchmark exists to make, and it
resolves into a single sentence: \textbf{India is the place this method is most
worth deploying and the place it is currently least trustworthy.}

\paragraph{Why India is the favourable case.} Delhi has the highest
cooling-degree-day share of any metered series in the benchmark at $0.89$, and
it is the only series whose demand rises \emph{strongly} with temperature
($+0.59$). France is its mirror at $-0.64$: French demand rises when it gets
cold, because the building stock is electrically heated. That distinction is the
whole economic argument. A controller that buys down a cooling peak inside a
demand-charge tariff has something to work with in the first regime and almost
nothing in the second. Delhi is also among the peakiest profiles in the set
($p99/$median $1.61$), which is precisely the shape a demand charge punishes.

\paragraph{Why India is the difficult case.} On both axes that determine whether
the controller can be relied upon, India is last among metered supplies. Skill
over a seasonal-naive baseline is $+0.097$, against $+0.713$ for Germany,
$+0.636$ for France and $+0.619$ for Spain --- a model that needs no training at
all is nearly as good on Delhi as ours. And coverage of the nominal $90\%$
interval is $0.762$ under the benchmark's calibration, the worst metered figure
anywhere in the study against $0.911$ for Ireland --- and $0.904$ once the
adaptive step is stated in the interval's units (Table~\ref{tab:aci-step}),
which is the reading that stands.

The second of these is the serious one. The controller's $q95$ is read from the
same object as that interval, so an interval that is too narrow yields a demand
ceiling that is too narrow, and the safety property is weakest on exactly the
load the project targets.

\begin{table}[t]
\centering
\footnotesize
\setlength{\tabcolsep}{3.5pt}
\resizebox{\linewidth}{!}{%
\begin{tabular}{llrrlrrrrr}
\toprule
Series & Level & June & Train mo. & Cooling on meter & Median kW & Seas.\ naive & Ours & Skill & Cov 90\% \\
\midrule
Campus & campus & 2017 & 8 & yes & 170.1 & 22.942 & 21.047 & \textbf{+0.083} & 0.873 \\
Academic & building & 2015 & 13 & no & 18.9 & 0.943 & 1.211 & \textbf{-0.284} & 0.903 \\
Academic & building & 2016 & 15 & no & 18.9 & 1.097 & 1.146 & \textbf{-0.045} & 0.910 \\
Academic & building & 2017 & 15 & no & 18.9 & 1.264 & 0.853 & \textbf{+0.325} & 0.909 \\
Boys & building & 2014 & 7 & no & 27.8 & 0.598 & 0.667 & \textbf{-0.114} & 0.915 \\
Dining & building & 2017 & 9 & no & 16.3 & 1.476 & 1.323 & \textbf{+0.103} & 0.901 \\
Facilities & building & 2015 & 10 & yes & 8.1 & 0.427 & 0.359 & \textbf{+0.160} & 0.902 \\
Facilities & building & 2016 & 5 & yes & 8.1 & 0.472 & 0.473 & \textbf{-0.001} & 0.897 \\
Facilities & building & 2017 & 15 & yes & 8.1 & 0.457 & 0.345 & \textbf{+0.246} & 0.903 \\
Girls & building & 2014 & 6 & no & 12.7 & 0.355 & 0.350 & \textbf{+0.016} & 0.905 \\
Girls & building & 2017 & 8 & no & 12.7 & 0.305 & 0.267 & \textbf{+0.124} & 0.908 \\
Library & building & 2014 & 7 & no & 6.2 & 0.546 & 0.553 & \textbf{-0.012} & 0.922 \\
Library & building & 2016 & 7 & no & 6.2 & 0.391 & 0.769 & \textbf{-0.965} & 0.941 \\
\bottomrule
\end{tabular}
}
\caption{The Indian building arm: I-BLEND, IIIT-Delhi, one row per admissible (series, June) window, native 15-minute blocks, the unchanged protocol with the training block shortened where meter uptime forces it (the China-arm precedent). \emph{Cooling on meter} is the selection rule's chilled-water clause: six of seven buildings are served by a central chiller on the campus feed, so only the Facilities building and the campus total are control objects; every row is a forecasting and calibration result. Over the 12 building rows the model beats seasonal naive on 6; mean skill is $-0.037$, $+0.060$ on the 4 windows with a full training block and $-0.086$ on the 8 shorter ones. Coverage is 0.897--0.941 on every building row and 0.873 on the campus feed.}
\label{tab:iblend}
\end{table}


\paragraph{At building level, the verdict is the same.} Table~\ref{tab:iblend}
is the Indian building arm: the IIIT-Delhi campus, thirteen windows, the
unchanged protocol. The skill column says what the city-level row said. Over
the twelve building windows the model beats a seasonal-naive baseline on six
and loses on six, with mean skill $-0.037$; the worst row in the entire study
is now an Indian one (Library, June 2016, $-0.965$, on a seven-month training
block). Skill tracks the length of the training block --- $+0.060$ on the four
windows with thirteen months or more, $-0.086$ on the eight shorter ones ---
which is the direction the China arm's adaptation was argued to err in, now
measured. The campus feed, the one object here that a demand charge is actually
levied on, returns $+0.083$: a model that needs no training is nearly as good.
That is India's skill verdict, and it no longer rests on a decade-old city
archive.

The calibration column says something the city-level row could not, and it
says it twice. Under the benchmark's calibration coverage holds at building
scale, $0.897$--$0.941$ on every window, exactly as it did on the eighteen BDG2
buildings; it is $0.873$ on the campus feed, and $0.762$ on the city.
Table~\ref{tab:calibration} carries this as its lower panel: one city, one
weather feed, three levels of aggregation, and coverage falling with each step
up the ladder --- which reproduced the first draft's aggregation effect inside
a single city with nothing else varying, and reproduced its confound with it.
Under the adaptive step stated in the interval's units the same three levels
read $0.908$, $0.888$ and $0.904$ (Table~\ref{tab:aci-step}): flat within
$0.02$, in a city where the step the building used was one sixth of the
building's on the campus feed and one twentieth on the city. The campus row is
the one that matters for control: the per-step exceedance of its $q95$ is
$0.079$ against a nominal $0.05$ under the benchmark's calibration, the only
window in Table~\ref{tab:horizon-panel} that is over the line. It is also the
one window in that table where the absolute step was \emph{slower} than the
relative one rather than faster ($\kappa=0.0015$ against $0.006$), which is
the campus feed being the largest Indian meter in the study; restated in the
interval's units its per-step rate falls to $0.069$ and its horizon exceedance
is unchanged at $0.438$ (Table~\ref{tab:horizon-kappa}).

\paragraph{What this arm cannot say.} Six of the seven buildings are on a
central chiller that is not on their meter, so the building rows are
forecasting and calibration results and no savings figure is claimed from them.
The campus feed is a control object --- it carries the chiller, and it is the
consumer the tariff bills --- and Section~\ref{sec:dial} runs the closed loop
on it under Delhi's own tariff, which closes the loop on an Indian consumer at
the resolution and under the tariff the bill is actually computed at. What that
run found, and did not find, is reported there.

\section{What India must fix}
\label{sec:india-fix}

Stated as work rather than as caveats, in descending order of how much they
cost.

\paragraph{1. No row of this study is both Indian and building-level.} The
Building Data Genome Project 2 harmonises $1{,}636$ buildings across three
countries and contains no Indian building, so the \emph{controller} was tested
on American buildings under an Indian tariff and the \emph{forecaster} on Indian
city demand, and no single row here is both. Relabelling an American building as
Indian would have been worse than the gap, so the gap is reported. An earlier
draft stated that no Indian building-level load dataset exists; that was wrong.
I-BLEND \citep{rashid2019iblend} provides 52 months (August 2013 to December
2017) of one-minute mains power and power factor for seven buildings on the
IIIT-Delhi campus --- academic, lecture, library, facilities, dining and two
dormitories --- under a CC0 licence. Three properties make it the right
corpus to close this gap and not merely a convenient one. It is in Delhi, the
city whose system demand is the study's focal series. It is sampled at one
minute, so the 15- and 30-minute billing blocks Indian demand charges are
assessed on can be formed by averaging rather than, as for every BDG2 building
here, interpolated up from hourly data that has already smoothed the peak. And
it spans four Junes, so the horizon-gap measurement of
Section~\ref{sec:horizon} can be replicated across years on one building rather
than reported once. What it lacks is submetering of HVAC, so the base-load
split remains a changepoint inference. The unchanged forecasting protocol has now been run on it
(Table~\ref{tab:iblend}); the closed loop under the Delhi (DERC) tariff has
not, and nothing in this paper anticipates its result.

\paragraph{2. Calibration collapses across the monsoon.} Coverage of $0.762$ is
an assumption failure, not a tuning failure: split conformal requires
calibration and test blocks to be exchangeable
\citep{vovk2005algorithmic,romano2019cqr}, and a Delhi May is not exchangeable
with a Delhi June. Section~\ref{sec:calibration} measures the remedy on this
feed: across a walk-forward year the adaptive layer, with its step carried over
from the building as a fraction of the interval width, holds $0.874$--$0.935$
by month where split conformal falls to $0.628$, and holds this June at
$0.895$; on the benchmark's own June, refitted and recalibrated
(Table~\ref{tab:aci-step}), it is $0.904$. The $0.762$ of
Table~\ref{tab:calibration} is the benchmark row under the step the building
used, and it is a units artefact rather than a monsoon one: the monsoon breaks
exchangeability, which split conformal cannot survive anywhere, and the layer
that survives it had been switched off by its units.

\paragraph{3. Indian load carries calendar structure the model does not encode.}
A large share of Indian public holidays are lunisolar --- Diwali, Holi, Eid ---
and move against the Gregorian calendar, so they cannot be expressed as the
fixed date list that load-forecasting pipelines conventionally use, including
ours until this study forced the change. Beyond holidays, Indian demand carries
load-shedding events, agricultural pumping schedules and festival periods that
have no counterpart in the European series and no representation in our feature
set. The low skill figure of $+0.097$ should be read partly as a statement about
missing features rather than solely about the data.

\paragraph{4. The Indian series is a decade old.} The Delhi archive covers
2011--2012. India's generation mix, air-conditioning penetration and peak demand
have all moved substantially since. The load \emph{shape} conclusions are the
durable ones; the levels are not, and nothing in this paper should be read as a
statement about the present Indian grid. We had a candidate answer to this
objection in the China arm and Section~\ref{sec:china} withdrew it.

\paragraph{What this does not undermine.} The favourable half of
Table~\ref{tab:india} rests on physical structure --- cooling share, the sign of
the load--temperature relationship, peakiness --- which is the part least likely
to have inverted since 2012. Indian demand has not stopped rising with
temperature.

\section{Results that argue against the method}
\label{sec:against}

Five rows are reported unchanged, and three of them are outright negative ---
worse than a baseline that needs no training at all.

The London office returns skill $-0.463$ and public services in Washington DC
returns $+0.052$. Both sit at $1$--$3\%$ HVAC share, where a weather-aware model
has little to add over ``last week, same time'' --- a coherent and testable
explanation rather than an excuse. Delhi returns $+0.097$ with the study's worst
metered calibration, and its explanation is different and given in
Section~\ref{sec:india-fix}: a system-level series carrying calendar structure
the feature set does not encode. Guangdong and Heilongjiang return $-0.314$ and
$-0.294$, the two lowest figures in the study, and Section~\ref{sec:china}
explains why those two should not be used against the method any more than they
should have been used for it.

Two defects in our own pipeline were exposed by running it outside its home
site. A hard-coded United States federal holiday calendar would have inverted
the holiday flag on two of the largest annual load anomalies for every European
building; holidays are now resolved per country, and India is the case that
forces it, its calendar being substantially lunisolar and therefore not
expressible as a fixed list. And a data-preparation stage overwrote its own
manifest, silently discarding every site but the last, which single-site work
had never exposed.

A third is worth recording because it is the kind that survives review. The
column commonly written as ``the Boole bound'' in applied chance-constraint work
is usually $1-(1-\alpha)^H$, which is the independence calculation and is not a
bound. Our own tables carried that label until this draft. The corrected reading
in Section~\ref{sec:horizon} is stronger than the one it replaces, because the
valid bound turns out to be vacuous over the horizon that matters.

\section{Limitations}

The system tier is a forecasting benchmark, not a control benchmark: there is no
thermal mass to shift at grid level and no savings claim is made from it. The
Delhi window (2011--12) does not overlap the European one (2016--17) or the
Chinese one (2018), which motivates reporting within-series skill rather than
absolute loss but does not eliminate the concern. Capital-city temperature
proxies national load-weighted temperature, acceptable for Great Britain and
Ireland and rough for Germany, France, Spain and every Chinese province.
Nothing here claims that model \emph{weights} transfer across countries --- our
own cold-start study finds the model losing to seasonal naive on an unseen
building at the same site, scoring $3.546$ pinball against that baseline's
$3.037$ --- only that the architecture and protocol do. The demographic arm is a
single city and cannot separate demographic from local building stock. Delhi is
one Indian city, and is real Indian data rather than a claim about India. The
China arm is reconstructed data and Section~\ref{sec:china} states exactly which
of its columns may be read. The acceptance result of Section~\ref{sec:dial} is
one building over one month; it is the deepest measurement in the study and the
narrowest.

The closed-loop results are produced on an RC plant model whose parameters are
fitted from each meter's own changepoint, but which holds a comfort band the
real building evidently does not. Simulated uncontrolled peak exceeds the
metered June peak by $9\%$ on the Phoenix office and $24\%$ on the Delhi
campus: the model air-conditions continuously where the real premises let
indoor temperature drift. Absolute savings from it are therefore an upper
bound, and we do not claim them as metered outcomes. The acceptance result is
unaffected in kind, because it grades the controller against the ceiling it
committed to rather than against a kilowatt figure, but a reader should take
the rupee numbers as the price of a risk level on this plant and not as a
realised bill.

The five arms were run once, with the adaptive step that
Section~\ref{sec:calibration} then found to be stated in the wrong units, and
they have deliberately not been rerun. Tables~\ref{tab:climate}--\ref{tab:china}
and Table~\ref{tab:calibration} are therefore reported as what that protocol
produced, and Table~\ref{tab:aci-step} --- the same forecaster, the same
supplies, the same predictions, calibrated three ways --- is the correction. We
prefer this to a silent regeneration for two reasons. A study that quietly
reruns itself after finding its own confound leaves a reader unable to see
either the confound or its size; and the pinball and skill columns of the five
arms do not depend on the question at issue, because every forecaster in a row
passes through the identical calibration layer, so the comparison between
forecasters is unaffected by which step that layer used. What the step changes
is the coverage column, and that column is re-reported in
Table~\ref{tab:aci-step} rather than overwritten.

\section{Conclusion}

We built an instrument rather than an algorithm and ran it unchanged over thirty
electricity supplies in eight countries to find out where a tariff-native demand
controller's guarantees actually hold.

The central result is that the guarantee most commonly deployed is not the
guarantee most commonly claimed. A per-step chance constraint on a demand
ceiling, marginally calibrated and behaving exactly as specified, is
under-protected over a sixteen-hour planning window by a factor of eight on the
building where it was first measured, and by between three and sixteen across
fourteen windows in two countries under either way of tuning the layer beneath
it; the
figure usually quoted as the conservative correction is not a bound and
overstates the risk by $2.4\times$; and the correction that is a bound is
vacuous four hours in. That risk level can be turned into an operator-set
parameter, at a measured price of \rupee{3{,}339} on a billing month and with
its inexactness and resolution floor stated. Alongside it, a finding this
paper first made and then unmade: distribution-free coverage does not degrade
with aggregation level. Split conformal fails at every tier alike, and the
adaptive layer holds the level at every tier once its step is stated in the
interval's units rather than the series' --- the ladder the first draft
reported was that parameter going inert as the meters grew. And forecast skill and available
flexibility are independent axes, only one of which tracks climate.

For India the benchmark returns a divided verdict, and both halves matter. India
has the most favourable load physics in the study --- the highest cooling share,
the only strongly positive load--temperature relationship, a peaky profile that
a demand-charge tariff punishes and this method is built to defend. It also
returns the lowest metered forecast skill and the worst metered calibration we
measured, and the loop has not yet been closed on an Indian building. The first
half is the case for the work. The second half is the work that remains, and
Section~\ref{sec:india-fix} states it as such rather than as a footnote ---
together with the Indian building corpus on which it can now be done.

Every number here is regenerated by one script from a clean checkout. The
results that contradict us are printed at the same size as the rest, and the one
result we most wanted to keep is the one Section~\ref{sec:china} withdraws.


\subsection*{Reproducibility}

Every number in this paper is written by a script from a file in
\texttt{results/}; none is typed into the text. From a clean checkout the whole
study regenerates in the order below: cache the system-level series and
I-BLEND, run the protocol over the thirty-seven supplies, then the Indian arm
window by window --- once with the adaptive step as an absolute number and once
as a fraction of interval width --- then the audits of the China arm, the
horizon bracket and closed loop, the calibration year, the step sweep, the
aggregation ladder under three calibrations, the copula's one free parameter,
the figures the prose states, and the references against their index. The last two commands emit every table
and figure. Each stage is resumable and skips work whose output already exists,
so an interrupted run continues rather than restarts.

{\footnotesize
\begin{verbatim}
python data/national.py
python data/iblend.py
python eval/comparative.py
./reproduce_india.sh
./reproduce_india_kappa.sh
python eval/china_audit.py
python eval/horizon_risk.py
python eval/conformal_audit.py
python eval/aci_gamma.py
python eval/aci_step_check.py
python eval/copula_df_check.py
python eval/paper_numbers.py
python eval/bibcheck.py
python eval/paper_tables.py
python eval/paper_figures.py
\end{verbatim}}

\noindent The calibration audit is run a second and third time on the Delhi
system series, with the protocol's own dates shifted five years onto that
series' archive, and the adaptive step stated first as the absolute number the
benchmark used and then as a fraction of interval width:

{\footnotesize
\begin{verbatim}
python eval/conformal_audit.py \
    --building IN_Delhi --shift-years 5
python eval/conformal_audit.py \
    --building IN_Delhi --shift-years 5 \
    --reuse --gamma-rel 0.006 --tag @rel
\end{verbatim}}

\noindent The second control object, the IIIT-Delhi campus feed, is billed
under the DERC schedule with no photovoltaics. Its demand target is found
first, and the closed loop is then run once at each adaptive step:

{\footnotesize
\begin{verbatim}
python eval/find_target.py \
    --building IIITD_Campus --tag @2017 \
    --tariff tariff/orders/derc_2021.json \
    --pv-kwp 0
python eval/horizon_risk.py \
    --building IIITD_Campus --tag @2017 \
    --tariff tariff/orders/derc_2021.json \
    --pv-kwp 0
python eval/horizon_risk.py \
    --building IIITD_Campus --tag @2017k \
    --tariff tariff/orders/derc_2021.json \
    --pv-kwp 0
\end{verbatim}}


% Statements Applied Energy requires of every submission.
\section*{Data availability}
Every dataset used is public. Building load is from the Building Data Genome
Project~2 \citep{miller2020bdg2} and from I-BLEND \citep{rashid2019iblend}, the
latter released under CC0. System-level demand is from the Delhi State Load Despatch Centre,
from ENTSO-E via Open Power System Data, and, for China, from the provincial
dataset of \citet{wu2023china}. The tariffs are the published DERC and TNERC
schedules cited in the text. The code that regenerates every table and figure
from these sources, in the order listed above, is at
\authortodo{repository URL, once public}.

\section*{CRediT authorship contribution statement}
% CRediT roles: Conceptualization, Methodology, Software, Validation, Formal
% analysis, Investigation, Resources, Data curation, Writing -- original draft,
% Writing -- review & editing, Visualization, Supervision, Project
% administration, Funding acquisition.
\noindent\textbf{Swetank~Kumar:} \authortodo{CRediT roles}.\\
\textbf{Devansh~Pathak:} \authortodo{CRediT roles}.\\
\textbf{Utkarsh~Sharma:} \authortodo{CRediT roles}.

\section*{Declaration of competing interest}
\authortodo{declare any relationships that could be seen to influence the work,
or confirm Elsevier's standard sentence: ``The authors declare that they have no
known competing financial interests or personal relationships that could have
appeared to influence the work reported in this paper.''}

\section*{Declaration of generative AI and AI-assisted technologies in the writing process}
\authortodo{Elsevier requires this where such tools were used, and they were.
Its template reads: ``During the preparation of this work the author(s) used
[tool] in order to [reason]. After using this tool, the author(s) reviewed and
edited the content as needed and take(s) full responsibility for the content of
the publication.'' Name the tool (Claude, Anthropic) and say what it was used
for, in your own words.}

\section*{Funding}
\authortodo{funding sources, or ``This research did not receive any specific
grant from funding agencies in the public, commercial, or not-for-profit
sectors.''}

\bibliographystyle{elsarticle-num-names}
\bibliography{../paper/refs}

\end{document}
